\maketitle
\noteprovenance

\begin{abstract}
\noindent
Let $p_0=2,p_1=3,\ldots$ enumerate the primes and let $g_n=p_{n+1}-p_n$.  Both
dyadic series converge and
\[
 \Pi=\sum_{n\ge0}\frac{p_n}{2^{\,n+1}}
 \;=\;2+\sum_{n\ge0}\frac{g_n}{2^{\,n+1}}\;=\;2+S .
\]
The identity is exact, carries no hypothesis, and rests on the elementary
polynomial bound $p_n\le1250(n+1)^4$ proved in Appendix~\ref{long251:app:prime-bound}.
The prime number theorem is not used.  Everything after it is negative
knowledge or an exact reformulation, and each statement is labelled by what it
is.

Write $T_N=\sum_{j\ge1}g_{N+j}2^{-j}$ for the scaled tails, so that
$T_{N+1}=2T_N-g_{N+1}$.  Irrationality of $\Pi$ is equivalent to
nonintegrality of every positive tail shift $T_{N+h}-T_N$, and equivalently to
the free-pair condition that for every modulus $t$ and every cutoff some pair
of indices congruent modulo $t$ beyond the cutoff has nonintegral tail
difference.  Two adjacent shifts in $(-1,1)$ with unequal corresponding gaps
cannot both be integral, and an explicit remainder bound certifies one such
pair at $h=1$, $N=2$ by exact integer arithmetic.  Every rational equal to
$\Pi$, and hence every rational equal to $S$, has denominator at least
$2^{589}>10^{177}$ by a certificate the Lean kernel decides, and at least
$2^{39997}>10^{12040}$ by a certified continued-fraction prefix.

Four proved statements constrain any argument.  Raising every gap by at most
$M$ beyond a prescribed prefix, preserving each residue modulo $M$, already
makes the sum rational.  That rationalising perturbation inherits the
fixed-block polynomial nonconcentration property which Schlage-Puchta proved
for the actual gaps, so nonconcentration together with prefix, size, residue
and cumulative-growth data still permits a rational value.  Under the two
window conditions the gap mismatch is forced to be $\pm2$, and by the same
nonconcentration lemma that event has density zero, so the missing input is a
cofinality statement on a sparse set.  Finally, two even values differing by
$2$, each recurring infinitely often inside every index residue class, do not
force irrationality: an explicit positive even word of prime-number-theorem
cumulative growth has both values recurring and an integral tail at every
index.  Problem~\#251 remains open, and the exact unproved input is cofinal
nonintegrality of the actual prime-gap tail shifts.
\end{abstract}

\begin{paperresult}{The exact identity, the finite exclusions, and the open boundary}
\textbf{Exact identity.} The prime-index dyadic series equals two plus the
consecutive-prime-gap dyadic series, unconditionally, with summability and the
identity checked by the Lean kernel over an elementary polynomial prime bound.
\textbf{Denominator floors.} Every rational equal to either series has
denominator at least $2^{589}>10^{177}$ by a kernel-decided certificate, and at
least $2^{39997}>10^{12040}$ by a certified continued-fraction prefix.
\textbf{Exact criteria.} Irrationality is equivalent to cofinal nonintegrality
of every fixed positive tail shift, to the free-pair condition with the offset
free, and to nonintegrality along the lcm diagonal.
\textbf{Proved obstructions.} Bounded residue-preserving perturbations make the
sum rational; that perturbation inherits fixed-block polynomial
nonconcentration; the two-window event forces a $\pm2$ mismatch and therefore
has density zero; and recurring gap values differing by $2$ inside every
residue class are compatible with an integral tail at every index.
\textbf{Open boundary.} No result here proves cofinal nonintegrality for the
actual prime gaps or irrationality of either series, and Erd\H{o}s~\#251
remains open.
\end{paperresult}

\section{Introduction}\label{long251:sec:problem}

Let $p_n$ denote the $n$th prime.  Erd\H{o}s Problem~\#251 asks whether
\[
 \Pi=\sum_{n\ge1}\frac{p_n}{2^{\,n}}
\]
is irrational~\cite[p.~94]{erdos1958}\cite[p.~62]{erdosgraham1980}%
\cite[p.~103]{erdos1988}.  Bloom's catalogue records the problem as
open~\cite{erdosproblems}.  Numerically
$\Pi=2/2+3/4+5/8+7/16+11/32+\cdots=3.674643966\ldots$.

Throughout the formal development the primes are indexed from zero, so that
$p^{(0)}_0=2$, $p^{(0)}_1=3$, $p^{(0)}_2=5$, and the gaps are
$g_i=p^{(0)}_{i+1}-p^{(0)}_i$.  In that indexing
\[
 \Pi=\sum_{i\ge0}\frac{p^{(0)}_i}{2^{\,i+1}} ,
\]
which is the convention every statement below uses; we drop the superscript
from here on.  The first values are
\[
 \begin{array}{c|cccccccccc}
  i   & 0&1&2&3&4 &5 &6 &7 &8 &9\\\hline
  p_i & 2&3&5&7&11&13&17&19&23&29\\
  g_i & 1&2&2&4&2 &4 &2 &4 &6 &2
 \end{array}
\]
so $g_0=1$ and every later gap is even, the primes after $2$ being odd.
A second normalisation with denominator $2^{\,i}$ also occurs,
and at every finite horizon it is exactly twice this one,
\[
 \sum_{i=0}^{n-1}\frac{p_i}{2^{\,i}}
 =2\sum_{i=0}^{n-1}\frac{p_i}{2^{\,i+1}} ,
\]
the
\lword{Erdos251/PrimeGapDyadicTail.lean}{111}{prime0DisplayedPartialSumQ_eq_two_mul}{factor-of-two
normalisation}.  A factor of two does not change rationality, so no result
below depends on the choice; the identity is stated so that the indexing cannot
drift silently.

\paragraph{Notation.}
$\N=\{0,1,2,\ldots\}$ and $\Npos=\{1,2,\ldots\}$.  We write $\ph$ for Euler's
totient function, $\operatorname{den}q$ for the denominator of a rational
number $q$ written in lowest terms, $\operatorname{dist}(x,\Z)$ for the
distance from a real number $x$ to the nearest integer, and
$\operatorname{Irr}(x)$ for the assertion that $x$ is irrational.  A rational
number is called \emph{integral} when it is the image of an integer.  A
sequence $(a_n)$ is \emph{eventually periodic with period $h\ge1$} when
$a_{n+h}=a_n$ for every sufficiently large $n$.  A set $A\subseteq\N$ has
\emph{density zero} when $|A\cap[1,x]|=o(x)$, and a property holds for
\emph{almost all} indices when its exceptional set has density zero.  A sum
over an empty range of indices is zero.

\paragraph{Relation to prior work.}
The passage from the primes to their gaps is already public.  Tao recorded on
the problem page that summation by parts makes the question equivalent to
irrationality of $\sum_n(p_{n+1}-p_n)2^{-n}$, and named the shape of the
missing input as a sufficiently quantitative and uniform prime-tuples
hypothesis giving statistical control of the binary expansion of about
$\log\log n$ consecutive gaps~\cite[comment of 7 October 2025]{erdosproblems251thread}.
Theorem~\ref{long251:res:infinite} below is that reduction in exact form, with the
endpoint retained, with convergence supplied by an elementary bound, and with
the statement checked by the Lean kernel.  The elementary bound replaces the
prime number theorem.  No priority is claimed for the reduction.

Land has since posted a conditional proof of irrationality assuming
Kuperberg's uniform Hardy-Littlewood prime-tuples conjecture, together with a
Lean formalisation of that conditional
argument~\cite[Conjecture~1 and \S1]{land2026}.  His route constructs weighted
prime-gap tails converging to $6$ from above, imposes a consecutive quadratic
prime pattern by a mixed prime-counting moment, and averages over later prime
indices to control the unprescribed tail.  Nothing in this note is
conditional on a prime-tuples hypothesis, and nothing in this note proves the
unconditional statement his argument would need.

Erd\H{o}s stated on p.~93 of the 1958 article that $\sum_n p_n^{k}/n!$ is
irrational for every $k\ge1$, and wrote there that the proof for $k>1$ is
complicated enough that only the case $k=1$ is
printed~\cite[pp.~93--95]{erdos1958}.  A proof of the full family appears in
Schlage-Puchta's Theorem~3, which gives the stronger statement that
$1,S_0,S_1,S_2,\ldots$ are linearly independent over $\Q$, where
$S_k=\sum_{n\ge1}p_n^{k}/n!$~\cite[Theorem~3]{schlagepuchta2011}.  The
catalogue entry still attributes the whole family to the 1958
article~\cite{erdosproblems}.

On page~103 of his 1988 problem paper Erd\H{o}s separately stated the
fixed-denominator problem: he could not prove that $\sum_n p_n^{k}/2^{n}$ is
irrational for every $k\ge1$, and wrote that the case $k=1$ was probably
already very difficult.  He also stated the variable-denominator expectation
that $\sum_{n\ge1}p_n/(g_1\cdots g_n)$ is irrational whenever $g_n\ge2$ and
$g_n=o(p_n)$~\cite[p.~103]{erdos1988}.  That expectation is false: Kova\v{c}
posted a construction of such a sequence $(g_n)$ for which the sum is exactly
$1$, on 15 April 2026~\cite[Theorem~1 and proof, pp.~1--2]{kovac2026}.  The
catalogue record still repeats the refuted expectation without mentioning the
counterexample, as checked on 6 September 2026~\cite{erdosproblems}; its
warning that relevant literature may be missing applies directly to this
omission.

Section~3 of the 1958 article proves a classification for the
variable-denominator family, and the hypotheses matter.  Let
$1<q_1\le q_2\le\cdots$ be integers satisfying, for some $k>0$, the lower
growth condition $q_n>o(n/\log^{k}n)$.  Then $\sum_{n\ge1}p_n/(q_1\cdots q_n)$
is rational if and only if $q_n=q\,p_n+1$ for one fixed integer $q\ge1$ and all
large $n$~\cite[pp.~96--97]{erdos1958}.  The denominators are nondecreasing,
and the printed condition bounds them from below.  Strict increase is not
assumed, and no upper bound is imposed.  The endpoint $q=1$ is attained: if
$G_n=\prod_{j=1}^{n}(p_j+1)$ then $p_n/G_n=1/G_{n-1}-1/G_n$, so the
corresponding series telescopes to $1$.  Kova\v{c}'s sequence evades the lower
growth condition, and neither statement addresses the fixed-denominator dyadic
series studied here.

There is also a genuinely adjacent proved dyadic theorem.  If $P(m)$ denotes
the largest prime factor of $m$, Erd\H{o}s and Pomerance proved that
\[
 \sum_{m\ge2}\frac{\mathbf 1_{\{P(m)>P(m+1)\}}}{2^m}
\]
is irrational~\cite[\S7, p.~320]{erdospomerance1978}.  Erd\H{o}s and Graham
record the complementary indicator on p.~62: equality of the two largest prime
factors is impossible for consecutive integers, so its series is $1/2$ minus
the displayed one and is irrational as well.  That is a theorem about a bounded
prime-factor comparison digit sequence, and the numerators in $\Pi$ are
unbounded.

Two further results of Schlage-Puchta enter below.  His Theorem~2 classifies
rationality for numbers whose base-$b$ digit string concatenates the
representations of a slowly growing integer sequence, and is a different
object from the tail criterion here; it was already pointed at this problem in
the catalogue thread on 15 April 2026.  His Lemma~4 is the input this note
actually uses.  It states that for every polynomial $F\in\Z[x_0,\ldots,x_k]$
which does not vanish identically, $F(g_n,\ldots,g_{n+k})\ne0$ for almost all
$n$, where $g_n$ are the actual consecutive prime
gaps~\cite[Lemma~4, pp.~5--6]{schlagepuchta2011}.  Its proof runs through
Selberg's sieve.  Sections~\ref{long251:sec:obstructions} draws two consequences from
it.

Problem~\#251 also appears as the unproved declaration \texttt{erdos\_251} in
the \emph{Formal Conjectures} repository~\cite{formalconjectures251}.  Its
zero-based Lean sum starts with the zeroth prime over $2^0$, so it is twice the
displayed normalisation and has equivalent irrationality status; its proof is
\texttt{sorry}.

\paragraph{The strategy.}
The digits $p_n$ grow, so the series is not a digit expansion in any bounded
alphabet, and the standard rationality criteria for such expansions do not
apply directly.  The classical elementary criteria for series of this kind
instead control irrationality through the growth of the denominators.
Erd\H{o}s and Straus named a sum $\sum_k 1/a_k$ over a strictly increasing
sequence of positive integers an \emph{Ahmes series}~\cite{erdosstraus1963};
for such a series the condition $a_k^{1/2^k}\to\infty$ is sufficient for
irrationality, and it is sharp, since shifted Sylvester sequences grow like
$C^{2^{k}}$ for arbitrarily large $C$ and have rational reciprocal sum.  Both
statements and their attribution are recorded in the introduction of Kova\v{c}
and Tao~\cite[\S1]{kovactao2024}.  Splitting each term $p_i/2^{\,i+1}$ into
$p_i$ copies of $2^{-(i+1)}$ writes $\Pi$ as a sum of unit fractions with
repetitions, and the denominators occurring in it are exactly the powers of
two.  Even ignoring the repetitions the growth hypothesis fails by every
available margin, since $(2^{\,n})^{1/2^{\,n}}\to1$, and every arithmetic
constraint has to come from the numerators instead.

What replaces growth control is denominator control on the sequence of
rescaled tails.  Rationality of the sum is equivalent to an eventual
integrality condition on differences of those tails, and that condition uses
nothing about the numerators beyond the fact that they are integers.  The
condition does not by itself force the numerators to repeat:
Proposition~\ref{long251:res:telescope}, applied to $K_n=n$, produces the integer
sequence $\kappa_n=n-1$, which is unbounded and hence not eventually periodic,
and whose dyadic sum is zero.

\paragraph{Outline.}
Section~\ref{long251:sec:parts} proves the identity and the irrationality equivalence.
Section~\ref{long251:sec:tail} develops the tail recurrence and states the three exact
criteria: pointwise shift escape, the free-pair form with a free offset, and
the lcm diagonal.  Section~\ref{long251:sec:local-certificate} gives the local
obstruction, its signed normal form, the explicit remainder bound and one
certified actual pair.  Section~\ref{long251:sec:denominator-floor} records the two
denominator floors.  Section~\ref{long251:sec:obstructions} proves the four
obstructions.  Section~\ref{long251:sec:measurements} reports the two finite
measurements.  Section~\ref{long251:sec:open} states the remaining obligation.  All
declarations linked below are checked by the Lean~4 kernel against the pinned
Mathlib revision at the formal-source commit~\commitshort, and the pinned
sources contain no proof placeholders and no project-defined axioms.  The
cited system paper identifies Lean~4~\cite[abstract and \S1,
pp.~625--626]{lean4}, while the Mathlib paper documents the library's
historical Lean~3-era architecture~\cite[abstract and \S1.1, p.~367]{mathlib};
it is not authority for the current pinned revision.

\statusboundary

\begingroup
\small
\begin{longtable}{@{}
  >{\raggedright\arraybackslash}p{0.34\linewidth}
  >{\raggedright\arraybackslash}p{0.18\linewidth}
  >{\raggedright\arraybackslash}p{0.40\linewidth}@{}}
\toprule
\papertablehead
Statement & Status & Treatment here \\
\midrule
Irrationality of $\Pi$ & Open & Not proved. \\
Finite prime-gap identity & Proved here & Theorem~\ref{long251:res:parts}, with the endpoint retained. \\
Infinite prime-gap identity & Lean-checked unconditionally & Theorem~\ref{long251:res:infinite}. \\
Prime-series/gap-series irrationality equivalence & Lean-checked unconditionally & Corollary~\ref{long251:res:irr-equivalence}. \\
Elementary polynomial prime bound $p_n\le1250(n+1)^4$ & Lean-checked; proof reprinted here & Appendix~\ref{long251:app:prime-bound}. \\
Actual real prime-gap tail and its recurrence & Lean-checked unconditionally & Section~\ref{long251:sec:tail}. \\
Rationality/integral-shift classification & Lean-checked; an equivalence & Theorem~\ref{long251:res:escape-irrational}. \\
Free-pair irrationality criterion & Lean-checked; an equivalence & Theorem~\ref{long251:res:freepair}. \\
Lcm-diagonal irrationality criterion & Lean-checked; an equivalence & Theorem~\ref{long251:res:lcmdiagonal}. \\
Adjacent small-shift obstruction, rational and real & Lean-checked & Theorem~\ref{long251:res:smallpair}, Corollary~\ref{long251:res:smallpair-real}. \\
Signed two-window normal form & Lean-checked; an equivalence & Theorem~\ref{long251:res:signedwindow}. \\
Explicit remainder bound and the certified pair at $h=1$, $N=2$ & Ordinary proof; exact integer replay in Appendix~\ref{long251:app:certificates} & Propositions~\ref{long251:res:explicit-remainder} and~\ref{long251:res:finite-smallpair}. \\
Denominator floor $b\ge2^{589}$ & Lean-checked by \texttt{decide +kernel}; ordinary proof reprinted here & Theorem~\ref{long251:res:denominatorfloor}. \\
Certified continued-fraction exclusion $q\ge2^{39997}$ & Exact computation over a Lean-checked tail bound & Theorem~\ref{long251:res:cfexclusion}; no Lean declaration. \\
Bounded-perturbation obstruction & Lean-checked & Theorem~\ref{long251:res:boundedperturbation}. \\
Nonconcentration survives the rationalising perturbation & Ordinary proof; external lemma cited & Theorem~\ref{long251:res:nonconcentration} and Corollary~\ref{long251:res:nonconc-primes}. \\
The two-window event has density zero & Ordinary proof; external lemma cited & Theorem~\ref{long251:res:sparse}. \\
Recurring gap values differing by $2$ do not suffice & Ordinary proof & Theorem~\ref{long251:res:polignacfail}. \\
Quadratic polynomial-shift countermodel & Lean-checked & Proposition~\ref{long251:res:polynomialcountermodel}. \\
Rationality alone forces periodic integer coefficients & False & Proposition~\ref{long251:res:telescope}. \\
Affine and fixed-lattice cofinal escape tests & Lean-checked; both equivalences & Theorem~\ref{long251:res:affinecollapse}. \\
Actual prime gaps are unbounded and not eventually periodic & Lean-checked & Proposition~\ref{long251:res:gap-nonperiodic}. \\
Adjacent-mismatch event density below $1.2\times10^{8}$ & Finite measurement, floating-point tails & Section~\ref{long251:sec:measurements}. \\
Free-pair witnesses for every $t\le20$ below $2\times10^{7}$ & Finite measurement & Section~\ref{long251:sec:measurements}. \\
Cofinal adjacent small mismatch & Sufficient condition; not proved & Problem~\ref{long251:prob:smallpair}. \\
\bottomrule
\end{longtable}
\endgroup

\smallskip
\noindent{\footnotesize
\emph{How to read the middle column.}  \emph{Proved here} marks a statement
proved in the text.  \emph{Lean-checked} marks a statement accepted by the
kernel at the pinned commit, with the declaration reachable through the linked
phrase.  \emph{Exact computation} marks a finite computation with a recorded
receipt and no Lean declaration.  \emph{Ordinary proof} marks a derivation
written out here with no formal check.  \emph{Finite measurement} marks a
computation over one bounded range which supplies no cofinal statement.\par}

\medskip
\noindent\textbf{Keywords.} irrationality; prime gaps; dyadic series;
summation by parts; Lean~4.
\qquad
\textbf{MSC 2020.} 11J72 (primary); 11N05, 68V20 (secondary).

\section{Summation by parts, with the endpoint retained}\label{long251:sec:parts}

Summation by parts trades a sequence for its consecutive differences.  The form
recorded here is exact: it carries no error term, it assumes nothing about the
sequence, and it retains the endpoint term.  No endpoint is absorbed into an
estimate.  For a sequence $P$ of rational numbers and $n\ge0$ put
\[
 D(P,n)=\sum_{i=0}^{n-1}\frac{P(i)}{2^{\,i+1}},
 \qquad
 \Delta(P,n)=\sum_{i=0}^{n-1}\frac{P(i+1)-P(i)}{2^{\,i+1}} ,
\]
both empty, hence zero, at $n=0$.  We call these the
\lword{Erdos251/PrimeGapDyadicTail.lean}{101}{dyadicPartialSumQ}{dyadic partial
sum} and the
\lword{Erdos251/PrimeGapDyadicTail.lean}{121}{dyadicDifferencePartialSumQ}{dyadic
difference sum} of $P$.

\begin{proposition}[finite summation by parts]\label{long251:res:abel}
For every rational sequence $P$ and every $n\ge0$,
\[
 D(P,n+1)=P(0)+\Delta(P,n)-\frac{P(n)}{2^{\,n+1}} .
\]
\end{proposition}

\begin{proof}
A routine induction on $n$.  At $n=0$ both sides equal $P(0)/2$.  For the step,
adding $P(n+1)/2^{\,n+2}$ to the left and $\bigl(P(n+1)-P(n)\bigr)/2^{\,n+1}$
to the difference sum changes the endpoint term from $P(n)/2^{\,n+1}$ to
$P(n+1)/2^{\,n+2}$, and the two adjustments agree.
\end{proof}

\noindent Formalised as the
\lword{Erdos251/PrimeGapDyadicTail.lean}{138}{dyadicPartialSumQ_eq_start_add_differences}{summation-by-parts
identity}.  Nothing is assumed about $P$: no positivity, no monotonicity, and
no convergence.

Specialising to $P(i)=p_i$, whose first value is $p_0=2$, and writing
$g_i=p_{i+1}-p_i$ for the zero-based gaps, gives the reformulation.

\begin{theorem}[prime-gap reformulation]\label{long251:res:parts}
Let $p_0=2,p_1=3,\ldots$ be the primes in increasing order and
$g_i=p_{i+1}-p_i$.  For every $n\ge0$,
\[
 \sum_{i=0}^{n}\frac{p_i}{2^{\,i+1}}
 =2+\sum_{i=0}^{n-1}\frac{g_i}{2^{\,i+1}}-\frac{p_n}{2^{\,n+1}} .
\]
\end{theorem}

\noindent Formalised as the
\lword{Erdos251/PrimeGapDyadicTail.lean}{172}{prime0_dyadic_summation_by_parts}{prime-gap
summation by parts}, using the
\lword{Erdos251/PrimeGapDyadicTail.lean}{161}{primeGapPartialSumQ}{gap partial
sum} and the
\lword{Erdos251/PrimeGapDyadicTail.lean}{156}{primeGap0_cast}{gap cast
identity}; the latter records that the natural-number difference
$p_{n+1}-p_n$ agrees with the difference taken in $\Q$, which needs
$p_n\le p_{n+1}$, the
\lword{Erdos251/PrimeGapDyadicTail.lean}{152}{prime0_mono_step}{monotonicity
step}.  The leading $2$ is the first prime, not a normalising constant.
At $n=2$, for instance, the left side is $2/2+3/4+5/8=19/8$ and the right side
is $2+(1/2+2/4)-5/8=19/8$.

\subsection{The infinite identity and the irrationality equivalence}\label{long251:sec:infinite}

Write
\[
 u_n=\frac{p_n}{2^{\,n+1}},\qquad
 v_n=\frac{g_n}{2^{\,n+1}}
\]
for the terms of the prime series and of the gap series, so that
$\sum_{n\ge0}u_n=\Pi$.  The termwise identity $v_n=2u_{n+1}-u_n$ is the
\lword{Erdos251/PrimeGapDyadicTail.lean}{202}{primeGapDyadicTerm_eq}{dyadic
discrete derivative}.  It expresses each gap term as an integer combination of
two consecutive prime terms, so once $(u_n)$ is summable the gap series can be
summed by rearranging two copies of the prime series.

\begin{theorem}[infinite prime-gap identity]\label{long251:res:infinite}
Both series converge and
\[
 \sum_{n\ge0}\frac{p_n}{2^{\,n+1}}
 \;=\;2+\sum_{n\ge0}\frac{g_n}{2^{\,n+1}} .
\]
\end{theorem}

\begin{proof}
The bound $p_n\le1250(n+1)^4$ of Appendix~\ref{long251:app:prime-bound} makes $(u_n)$
summable, and $0<g_n\le p_{n+1}$ then makes $(v_n)$ summable.  The shifted
sequence $(u_{n+1})$ is summable, so summing $v_n=2u_{n+1}-u_n$ and using
$u_0=1$ gives
\[
 \sum_{n\ge0}v_n
 =2\left(\sum_{n\ge0}u_n-u_0\right)-\sum_{n\ge0}u_n
 =\sum_{n\ge0}u_n-2 . \qedhere
\]
\end{proof}

\noindent The polynomial bound is
\lword{Erdos251/PrimeGapDyadicTail.lean}{360}{prime0_le_polynomial}{checked
here}, summability of the prime series
\lword{Erdos251/PrimeGapDyadicTail.lean}{379}{summable_primeDyadicTerm}{here},
the summability transfer is the
\lword{Erdos251/PrimeGapDyadicTail.lean}{385}{summable_primeGapDyadicTerm_of_summable_primeDyadicTerm}{gap-series
summability theorem}, and the displayed identity is the
\lword{Erdos251/PrimeGapDyadicTail.lean}{427}{tsum_primeDyadicTerm_eq_two_add_primeGap_unconditional}{unconditional
prime-gap identity}.  The prime number theorem is not a dependency at any
point.

\begin{corollary}[exact irrationality reformulation]\label{long251:res:irr-equivalence}
$\Pi$ is irrational if and only if $S=\sum_{n\ge0}g_n2^{-(n+1)}$ is irrational.
The corresponding zero-based series with denominator $2^n$ equals $4+2S$ and
has the same irrationality status.
\end{corollary}

\noindent These are the
\lword{Erdos251/PrimeGapDyadicTail.lean}{435}{irrational_tsum_primeDyadicTerm_iff_primeGap}{normalised
irrationality equivalence}, the
\lword{Erdos251/PrimeGapDyadicTail.lean}{444}{tsum_primeDisplayedDyadicTerm_eq_four_add_two_primeGap}{displayed-series
identity}, and the
\lword{Erdos251/PrimeGapDyadicTail.lean}{459}{irrational_tsum_primeDisplayedDyadicTerm_iff_primeGap}{displayed-series
irrationality equivalence}.  Concretely, $\Pi=3.674643966\ldots$ is irrational
if and only if $S=1.674643966\ldots$ is, and neither is known.

\section{The tail recurrence and the exact criteria}\label{long251:sec:tail}

The series is not attacked directly.  Suppose $\sum_{i\ge0}a_i2^{-(i+1)}$
converges with every $a_i$ an integer, and rescale its tails by putting
\[
 T_N=2^{\,N+1}\sum_{i>N}\frac{a_i}{2^{\,i+1}}
    =\sum_{j\ge1}\frac{a_{N+j}}{2^{\,j}} .
\]
Two facts follow immediately.  First $T_{N+1}=2T_N-a_{N+1}$, so moving one
level along doubles the rescaled tail and subtracts a single coefficient.
Second $T_0=2\sum_{i\ge0}a_i2^{-(i+1)}-a_0$, so the sum is rational exactly
when $T_0$ is.  This section studies that recurrence, assuming nothing about
the coefficients beyond the fact that they are integers, and resumes the
prime-gap instance at the end.

\begin{definition}\label{long251:def:rec}
Let $g:\N\to\Z$ and $T:\N\to\Q$ or $T:\N\to\R$.  Say $T$ satisfies the
\emph{dyadic tail recurrence} with \emph{digits} $g$ when
$T_{N+1}=2T_N-g_{N+1}$ for every $N$.  Write $\sigma_h(N)=T_{N+h}-T_N$ for the
\emph{shift} of length $h$ at $N$, and call a rational number \emph{integral}
when it is the image of an integer.
\end{definition}

\noindent These are the
\lword{Erdos251/PrimeGapDyadicTail.lean}{482}{DyadicTailRecurrence}{tail
recurrence}, the
\lword{Erdos251/PrimeGapDyadicTail.lean}{544}{tailShift}{shift}, and
\lword{Erdos251/PrimeGapDyadicTail.lean}{575}{RatIntegral}{integrality}.
One small orbit is worth having in view.  Take every digit $g_N=0$ and
$T_0=1/12$.  Then the orbit is
$\tfrac1{12},\tfrac16,\tfrac13,\tfrac23,\tfrac43,\ldots$ and
$\sigma_h(N)=(2^{h}-1)2^{\,N}/12$.  The shift of length $2$ is not integral at
$N=0$ and is integral at $N=2$; the shift of length $1$ is never integral.
Nonintegrality at one prescribed shift length therefore does not imply
irrationality.

Iterating the recurrence $h$ times multiplies $T_N$ by $2^{h}$ and accumulates
an explicit integer.  Define $B_{0,N}=0$ and $B_{h+1,N}=2B_{h,N}+g_{N+h+1}$,
so that $B_{h,N}=g_{N+1}2^{\,h-1}+\cdots+g_{N+h}$: the
\lword{Erdos251/PrimeGapDyadicTail.lean}{647}{dyadicTailBlock}{tail block}.

\begin{theorem}[block identity]\label{long251:res:block}
For every $N$ and $h$,
\[
 T_{N+h}=2^{h}T_N-B_{h,N},
 \qquad\text{hence}\qquad
 \sigma_h(N)=(2^{h}-1)\,T_N-B_{h,N} .
\]
\end{theorem}

\begin{proof}
A routine induction on $h$; the step is one application of the recurrence
together with the recursion defining $B$.  The second identity is the first
minus $T_N$.
\end{proof}

\noindent Formalised as the
\lword{Erdos251/PrimeGapDyadicTail.lean}{653}{tail_iterate_eq_pow_mul_sub_block}{iterated
block identity} and the
\lword{Erdos251/PrimeGapDyadicTail.lean}{667}{tailShift_eq_scaled_sub_block}{scaled
shift identity}.  The shift also obeys the recurrence in its own right,
\begin{equation}\label{long251:eq:shift-step}
 \sigma_h(N+1)=2\sigma_h(N)-\bigl(g_{N+h+1}-g_{N+1}\bigr),
\end{equation}
the
\lword{Erdos251/PrimeGapDyadicTail.lean}{563}{tailShift_succ}{shift step
identity}.  Since $B_{h,N}$ is an integer, Theorem~\ref{long251:res:block} converts a
question about the shift into a question about $(2^{h}-1)T_N$: the shift
$\sigma_h(N)$ is integral if and only if $(2^{h}-1)T_N$ is, the
\lword{Erdos251/PrimeGapDyadicTail.lean}{802}{tailShift_integral_iff_scaledTail}{integral-shift
criterion}.  The criterion holds in both directions, so once
$T_N$ is fixed no choice of the digits after index $N$ can change whether
$\sigma_h(N)$ is an integer.

For a rational orbit the denominator settles the matter exactly:
\begin{equation}
 \operatorname{den}(T_{N+1})
 =\frac{\operatorname{den}(T_N)}{\gcd(2,\operatorname{den}(T_N))},
 \qquad
 \sigma_h(N)\in\Z
 \iff
 \operatorname{den}(T_N)\mid 2^h-1 .
\tag{3.1}\label{long251:eq:den-law}
\end{equation}
Each even denominator loses exactly one factor of $2$ and an odd denominator is
unchanged, so after finitely many steps the denominator is an odd integer $d$;
Euler's congruence then gives $d\mid2^{\ph(d)}-1$, and the multiplicative order
of $2$ modulo $d$ is the least such exponent.  These are the
\lword{Erdos251/PrimeGapDyadicTail.lean}{1217}{tail_den_succ}{denominator
recurrence}, its
\lword{Erdos251/PrimeGapDyadicTail.lean}{1226}{tail_den_succ_eq_of_odd}{odd
case} and
\lword{Erdos251/PrimeGapDyadicTail.lean}{1236}{tail_den_succ_eq_half_of_even}{even
case}, the
\lword{Erdos251/PrimeGapDyadicTail.lean}{1279}{tailShift_integral_iff_den_dvd_mersenne}{denominator
classification}, the
\lword{Erdos251/PrimeGapDyadicTail.lean}{877}{tailShift_integral_totient_of_odd_den}{totient
shift}, and the
\lword{Erdos251/PrimeGapDyadicTail.lean}{900}{tailShift_integral_add}{propagation
of an integral shift to every later index}.  The hypothesis that $d$ is odd
cannot be dropped: if $T_N=1/2$ then $2^{h}-1$ is odd for every $h\ge1$, so no
shift at $N$ is integral.

\begin{theorem}[exact rationality classification]\label{long251:res:escape-irrational}
Let $T:\N\to\R$ satisfy $T_{N+1}=2T_N-g_{N+1}$ with integer digits $g$.  The
following are equivalent:
\begin{enumerate}[label=\textup{(\roman*)}]
\item $T_0$ is rational;
\item $\sigma_h(N)$ is an integer for some $h\ge1$ and some $N$;
\item for some fixed $h\ge1$, $\sigma_h(N)$ is an integer at every
  sufficiently large $N$.
\end{enumerate}
Consequently $T_0$ is irrational if and only if every positive-length shift is
nonintegral at every index, equivalently if and only if for every $h\ge1$ and
every cutoff some later $N$ has $\sigma_h(N)\notin\Z$.
\end{theorem}

\begin{proof}
For \textup{(i)}$\Rightarrow$\textup{(iii)}, the block identity identifies the
whole orbit with the cast of the rational recurrence starting at $T_0$; apply
\eqref{long251:eq:den-law} at an index where the denominator has become odd, and
propagate forward.  The implication
\textup{(iii)}$\Rightarrow$\textup{(ii)} is immediate.  For
\textup{(ii)}$\Rightarrow$\textup{(i)}, the block identity gives
$\sigma_h(N)=(2^h-1)T_N-B_{h,N}$ with $B_{h,N}$ and $\sigma_h(N)$ integers and
$2^h-1\ne0$, so $T_N$ is rational, and $T_N=2^NT_0-B_{N,0}$ then makes $T_0$
rational.  Negating the three conditions gives the two irrationality
formulations.
\end{proof}

\noindent The exact real classifiers are
\lword{Erdos251/PrimeGapDyadicTail.lean}{1479}{not_irrational_initial_iff_exists_integral_positive_tailShift}{one
integral positive shift},
\lword{Erdos251/PrimeGapDyadicTail.lean}{1525}{not_irrational_initial_iff_exists_eventually_integral_positive_tailShift}{eventual
integrality},
\lword{Erdos251/PrimeGapDyadicTail.lean}{1551}{irrational_initial_iff_all_positive_tailShifts_nonintegral}{pointwise
nonintegrality}, and
\lword{Erdos251/PrimeGapDyadicTail.lean}{1572}{irrational_initial_iff_cofinalNonintegralTailShifts}{cofinal
nonintegrality}.

\paragraph{The actual prime-gap orbit.}
Put
\[
 T_N=\sum_{j\ge1}\frac{g_{N+j}}{2^{\,j}},
\]
which converges by Theorem~\ref{long251:res:infinite}.  The scaled tail equals its
shifted-gap series, the
\lword{Erdos251/RealPrimeGapTail.lean}{31}{realPrimeGapTail_eq_tsum_shifted_gaps}{shifted-gap
series identity}, and satisfies $T_{N+1}=2T_N-g_{N+1}$ with no rationality
hypothesis, the
\lword{Erdos251/RealPrimeGapTail.lean}{48}{realPrimeGapTail_recurrence}{real
tail recurrence}.  Since $T_0=2S-1=2.349287932\ldots$, the numbers $\Pi$, $S$
and $T_0$ have the same rationality status, and
Theorem~\ref{long251:res:escape-irrational} specialises to the actual series:
irrationality of $\Pi$ is equivalent to cofinal nonintegrality of the positive
shifts of $T$, the
\lword{Erdos251/RealPrimeGapTail.lean}{73}{irrational_primeSeries_iff_realPrimeGapTail_escape}{escape
equivalence}.  Nothing about the bridge from the series to its actual tail is
left open; the open content is the escape hypothesis itself.

\subsection{The free-pair criterion and the lcm diagonal}\label{long251:sec:freepair}

The escape condition quantifies over a fixed shift length.  Freeing the offset
gives an equivalent condition with a weaker-looking producer.

\begin{theorem}[free-pair criterion]\label{long251:res:freepair}
$S$ is irrational if and only if for every $t\ge1$ and every $N_0$ there are
$N,M\ge N_0$ with $M\equiv N\pmod t$ and $T_M-T_N\notin\Z$.
\end{theorem}

\begin{proof}
If $S$ is irrational, take $M=N+t$ for the $N$ supplied by
Theorem~\ref{long251:res:escape-irrational} at shift length $t$.  Conversely, suppose
$S$ is rational.  By \eqref{long251:eq:den-law} the orbit reaches an index $N_0$ beyond
which the reduced denominator is a fixed odd $d$; let $t$ be the
multiplicative order of $2$ modulo $d$.  For $N,M\ge N_0$ with $M\equiv N$
modulo $t$, the difference $T_M-T_N$ is then an integer, contradicting the
stated property at this $t$ and this cutoff.
\end{proof}

\noindent The offset $M-N$ is free, so one nonintegral congruent pair for each
modulus and each cutoff suffices in place of a supply at a single fixed offset.
The mechanism is exact: beyond a rational state with odd reduced denominator
$d$, the difference $T_M-T_N$ is integral precisely when $M\equiv N$ modulo the
order of $2$ in $\Z/d\Z$, the
\lword{Erdos251/FreePairReduction.lean}{92}{free_pair_integral_iff_modEq}{free-pair
lattice}, so every rational orbit has a cutoff and a modulus at which
integrality holds exactly on the congruence classes, the
\lword{Erdos251/FreePairReduction.lean}{105}{exists_free_pair_lattice}{lattice
corollary}.  The criterion for the actual series is checked as the
\lword{Erdos251/FreePairReduction.lean}{236}{irrational_primeGap_tsum_iff_cofinalFreePairNonintegral}{free-pair
equivalence}, over the
\lword{Erdos251/FreePairReduction.lean}{186}{primeGapRealTail_recurrence}{actual
real tail recurrence}.  No antecedent for this packaging has been located; the
nearest located neighbour is Schlage-Puchta's
Theorem~2~\cite{schlagepuchta2011}, which treats periodic base-$b$ digit
expansions and is a different object, and the sweep of the dyadic
tail-recurrence and B\'ezivin-style functional-equation literature has not been
run.

The criterion rules out two shortcuts.  Rationality makes the tail a
finite-state object at every scale, so pairs with $T_M=T_N$ are abundant for
free and the whole difficulty sits in nonintegrality of the difference.  And
under the free-pair lattice the difference is an integer whenever the modulus
divides the offset, and an integer cannot lie in $(\tfrac12,1)$, so the window
component of the reduced event below is already a contradiction on its own.
Both derivations are ordinary proofs, recorded in the long record and not
formalised.

A second exact reformulation collapses all shift lengths and basepoints onto a
single schedule.  Let $L_0=1$ and $L_j=\lcm(1,\ldots,j)$.

\begin{theorem}[\lword{Erdos251/OrderLatticeDiagonal.lean}{153}{irrational_initial_iff_all_lcmDiagonal_nonintegral}{lcm-diagonal criterion}]\label{long251:res:lcmdiagonal}
Let $g:\N\to\Z$ and $T:\N\to\R$ satisfy $T_{N+1}=2T_N-g_{N+1}$.  Then
$T_0$ is irrational if and only if
$T_{2L_j}-T_{L_j}\notin\Z$ for every $j\ge0$.
\end{theorem}

\begin{proof}
Theorem~\ref{long251:res:escape-irrational} gives the forward implication.  Conversely,
if $T_0$ is rational then some positive shift length $h$ is integral at every
basepoint beyond an index $N_0$.  Choose $j$ so large that $N_0\le L_j$ and
$h\mid L_j$.  The cocycle identity
$\sigma_{a+b}(N)=\sigma_a(N)+\sigma_b(N+a)$ shows by induction that every
positive multiple of $h$ is integral at every such basepoint, so
$T_{2L_j}-T_{L_j}$ is an integer.
\end{proof}

The factorial schedule has the same divisibility property; the lcm schedule is
pointwise no larger and is the endpoint used here.  Theorems~\ref{long251:res:freepair}
and~\ref{long251:res:lcmdiagonal} are exact reformulations on the class of
integer-digit orbits.  They change the shape of the required producer and they
supply no nonintegrality statement for the actual prime gaps.

\section{A local certificate, and one actual pair}\label{long251:sec:local-certificate}

The classifiers reduce irrationality to a condition of infinite precision
imposed on a complete tail.  Inside $(-1,1)$ integrality is equality with zero,
so on that range the shift step identity \eqref{long251:eq:shift-step} turns
simultaneous integrality of two adjacent shifts into a single comparison of
digits.  What this buys is a certificate attached to one pair of adjacent
indices, whose verification already contradicts integrality there.  The
distance of a shift from the integers never has to be estimated.

\begin{theorem}[adjacent small-shift obstruction]\label{long251:res:smallpair}
Let $T$ satisfy the dyadic tail recurrence with integer digits $g$, and fix
$h$ and $N$.  If
\[
 -1<\sigma_h(N)<1,\qquad -1<\sigma_h(N+1)<1,
 \qquad g_{N+h+1}\ne g_{N+1},
\]
then $\sigma_h(N)$ and $\sigma_h(N+1)$ cannot both be integers.  Consequently,
if such a pair occurs beyond every threshold, the $h$-shift is not eventually
integral.
\end{theorem}

\begin{proof}
An integer strictly between $-1$ and $1$ is zero.  If both shifts were
integers, both would vanish, and \eqref{long251:eq:shift-step} would give
$g_{N+h+1}=g_{N+1}$.  The cofinal statement chooses one such adjacent pair
after the alleged onset of integrality.
\end{proof}

\begin{corollary}[real form and the sufficient condition]\label{long251:res:smallpair-real}
The same statement holds for a real orbit, with the same proof.  If for every
$h\ge1$ and every cutoff some later $N$ satisfies the three displayed
conditions for the actual prime gaps, then $\Pi$ is irrational.
\end{corollary}

\noindent The rational finite contradiction is the
\lword{Erdos251/PrimeGapDyadicTail.lean}{979}{tailShift_not_both_integral_of_small_pair_of_digit_ne}{adjacent
small-shift obstruction} with its
\lword{Erdos251/PrimeGapDyadicTail.lean}{1006}{not_eventuallyIntegralTailShift_of_cofinal_small_mismatch}{cofinal
form}; the real form is the
\lword{Erdos251/RealPrimeGapTail.lean}{99}{realTailShift_not_both_integral_of_small_pair_of_digit_ne}{real
adjacent-pair consumer}, and the implication to irrationality of the actual
series is the
\lword{Erdos251/RealPrimeGapTail.lean}{120}{irrational_primeSeries_of_realPrimeGapTail_small_mismatch}{small-mismatch
criterion}.  The third hypothesis alone is available for the actual gaps: it
reads $g_4=2\ne4=g_3$ at $h=1$, $N=2$, and by
Proposition~\ref{long251:res:gap-nonperiodic} it holds for arbitrarily large $N$ at
every fixed $h$.  What is missing is the pair of inequalities, each of which
constrains a complete infinite tail.

\begin{proposition}[prime gaps do not become periodic]\label{long251:res:gap-nonperiodic}
For every positive $h$, the actual consecutive-prime-gap sequence is not
eventually periodic with period $h$.
\end{proposition}

\begin{proof}
For $m\ge2$ the integers $(m+1)!+2,\ldots,(m+1)!+m+1$ are composite, so the
gaps are unbounded.  An eventually periodic sequence of natural numbers has
finite range after its preperiod and a bounded initial segment, hence is
bounded.
\end{proof}

\noindent Lean checks the factorial argument as
\lword{Erdos251/PrimeGapDyadicTail.lean}{57}{exists_primeGap0_gt}{unboundedness
of the actual gaps} and the conclusion as
\lword{Erdos251/PrimeGapDyadicTail.lean}{1023}{primeGap0_not_eventually_periodic}{non-eventual
periodicity}.  Far stronger lower bounds for large gaps are
known~\cite[Theorem~1, p.~66]{fgkmt2018}; unboundedness is all that is needed.

The conjunction in Theorem~\ref{long251:res:smallpair} has an exact normal form.  For
fixed $h$ write $D_N=\sigma_h(N)$ and $\delta_N=g_{N+h+1}-g_{N+1}$, so that
$D_{N+1}=2D_N-\delta_N$.

\begin{theorem}[\lword{Erdos251/AffineShiftEscape.lean}{113}{adjacent_small_mismatch_iff_signed_two_window}{signed two-window normal form}]\label{long251:res:signedwindow}
Assume $\delta_N$ is even.  The conjunction
$-1<D_N<1$, $-1<D_{N+1}<1$, $\delta_N\ne0$ is equivalent to
\[
 \bigl(\delta_N=2\ \hbox{ and }\tfrac12<D_N<1\bigr)
 \quad\hbox{or}\quad
 \bigl(\delta_N=-2\ \hbox{ and }-1<D_N<-\tfrac12\bigr).
\]
\end{theorem}

\begin{proof}
The recurrence and the two unit windows give $-3<\delta_N<3$.  A nonzero even
integer in that interval is $2$ or $-2$.  Substitution into
$D_{N+1}=2D_N-\delta_N$ gives the stated half-window and, in the reverse
direction, recovers the second unit window.
\end{proof}

For the actual gaps and $N\ge1$ every $\delta_N$ is even, so
Theorem~\ref{long251:res:signedwindow} applies throughout the range where the local
certificate is sought.  The same computation for a real orbit is one line and
is used in Theorem~\ref{long251:res:sparse} below.

\subsection{An explicit remainder and a certified actual pair}

Both window conditions involve complete infinite tails.  A dominated
truncation reduces each of them to a finite integer comparison.

\begin{proposition}[explicit remainder]\label{long251:res:explicit-remainder}
For integers $h,N\ge0$ and $L\ge1$ put
\[
 F_{h,N,L}=\sum_{j=1}^{L}\frac{g_{N+h+j}-g_{N+j}}{2^{\,j}},\qquad
 P(x)=x^4+8x^3+36x^2+104x+150,
\]
\[
 E_{h,N,L}=\frac{1250}{2^{L}}\bigl(P(N+h+L+2)+P(N+L+2)\bigr).
\]
Then $\bigl|\sigma_h(N)-F_{h,N,L}\bigr|\le E_{h,N,L}$.  In particular
$|F_{h,N,L}|+E_{h,N,L}<1$ certifies $|\sigma_h(N)|<1$, and
$\operatorname{dist}(F_{h,N,L},\Z)>E_{h,N,L}$ certifies
$\sigma_h(N)\notin\Z$.
\end{proposition}

\begin{proof}
The checked bound $p_n\le1250(n+1)^4$ gives $0\le g_m\le p_{m+1}\le1250(m+2)^4$,
so the omitted absolute tail is at most
\[
 1250\sum_{j>L}\frac{(N+h+j+2)^4+(N+j+2)^4}{2^{\,j}} .
\]
The polynomial identity $2P(x)=(x+1)^4+P(x+1)$ telescopes to
$\sum_{k=1}^{J}(m+k)^42^{-k}=P(m)-P(m+J)2^{-J}$, whose last term tends to zero.
Apply it with $m=N+h+L+2$ and $m=N+L+2$ after writing $j=L+k$.  The two tests
follow from the triangle inequality and the definition of distance to $\Z$.
\end{proof}

\begin{proposition}[a certified adjacent pair]\label{long251:res:finite-smallpair}
For the actual prime gaps, $h=1$ and $N=2$ satisfy the three hypotheses of
Theorem~\ref{long251:res:smallpair}: both $\sigma_1(2)$ and $\sigma_1(3)$ lie in
$(-1,1)$ and are nonintegral, and $g_4=2\ne4=g_3$.
\end{proposition}

\begin{proof}
Take $L=40$ and $Q=2^{40}=1099511627776$.  Exact integer arithmetic over the
first $46$ primes gives
\[
\begin{array}{c|r|r}
 N&Q\,F_{1,N,40}&Q\,E_{1,N,40}\\\hline
 2&-662838684750&11764181250\\
 3& 873345886050&12805761250
\end{array}
\]
In each row $|QF|+QE<Q$, which puts the whole certified interval inside
$(-1,1)$, and $QE$ is smaller than the distance from $QF$ to the nearest
multiple of $Q$, which excludes every integer.  The margins are $424908761776$
and $213359980476$ respectively.  These are strict integer comparisons rather
than inferences from rounded numerical tails.
Appendix~\ref{long251:app:certificates} replays them.
\end{proof}

This witness proves the local condition at one pair of indices.  The
quantifiers over every $h$ and every cutoff remain.

\section{Two denominator floors}\label{long251:sec:denominator-floor}

The same polynomial bound yields an exact rational enclosure of $\Pi$, and any
such enclosure excludes an initial range of denominators.

\begin{theorem}[kernel-decided denominator floor]\label{long251:res:denominatorfloor}
Let $a\in\Z$ and $b\in\Npos$.  If $\Pi=a/b$ then $b\ge2^{589}>10^{177}$, and
the same floor holds for every rational equal to $S$.
\end{theorem}

\begin{proof}
Put $c=1229$, the number of primes below $10^4$, and
$A=\sum_{i=0}^{c-1}p_i2^{\,c-i-1}$, so that $A/2^{c}$ is the $c$th partial sum
of $\Pi$.  All terms are positive, so $A/2^{c}\le\Pi$.  For the upper bound,
$p_{c+j}\le1250(c+j+1)^4$ and $(c+1+j)^4\le(c+1)^4(3/2)^{j}$ for $j\ge0$, valid
because $c\ge9$; summing the resulting geometric tail with ratio $3/4$ gives
\[
 \frac{A}{2^{c}}\;\le\;\Pi\;\le\;\frac{2A+5000(c+1)^4}{2^{\,c+1}} .
\]
The four positive integers $u,v,u',v'$ printed in
Appendix~\ref{long251:app:certificates} satisfy
\[
 u'v-uv'=1,\qquad
 \frac uv<\frac{A}{2^{c}},\qquad
 \frac{2A+5000(c+1)^4}{2^{\,c+1}}<\frac{u'}{v'},\qquad
 v+v'\ge2^{589},
\]
all four being integer comparisons after clearing the positive denominators.
If $\Pi=a/b$ then $av-bu\ge1$ and $bu'-av'\ge1$, and the determinant identity
gives
\[
 b=b(u'v-uv')=(av-bu)v'+(bu'-av')v\ge v+v' \ge 2^{589}.
\]
Since $589\log_{10}2>177$, the floor exceeds $10^{177}$.  If $S=a/b$, apply the
same argument to $\Pi=(a+2b)/b$.
\end{proof}

\noindent The Lean kernel decides the four inequalities on the same literals by
\texttt{decide +kernel} with no \texttt{native\_decide}, over a trial-division
sieve it re-runs on $[2,10^4)$: the
\lword{Erdos251/KernelDenominatorFloor.lean}{310}{cert_10000}{certificate},
the
\lword{Erdos251/KernelDenominatorFloor.lean}{316}{kernel_denominator_floor}{floor
for the prime series}, and the
\lword{Erdos251/KernelDenominatorFloor.lean}{324}{kernel_denominator_floor_primeGap}{floor
for the gap series}.  Its analytic input is the same enclosure, over the
polynomial bound at
\lword{Erdos251/PrimeGapDyadicTail.lean}{360}{prime0_le_polynomial}{line~360}.
The four Farey literals have $175$ to $181$ digits.

\begin{theorem}[certified continued-fraction exclusion]\label{long251:res:cfexclusion}
Every rational equal to $\Pi$, and hence every rational equal to $S$, has
reduced denominator $q\ge2^{39997}$, and therefore $q>10^{12040}$.
\end{theorem}

\emph{Evidence.}  This is an exact finite computation and no Lean declaration
carries it.  The continued-fraction expansion of $\Pi$ is developed to $23369$
partial quotients at a working scale of $80000$ bits.  Each quotient is forced
by a rational bracket of width $1$ in those scaled integer units around the
true value, and the separation $|q_n\Pi-p_n|>0$ is verified against that
bracket, so the prefix is the true prefix.  No floating-point reconstruction
enters it.  The bracket uses the Lean-checked tail estimate
$p_n\le1250(n+1)^4$, so no conjecture enters.  The classical theorem that a
best approximation of the second kind is a convergent then converts the
certified prefix into the denominator bound~\cite{khinchin1964}.  The receipt
is \texttt{state/formal\_math/probes/erdos251\_certified\_cf\_receipt.json},
which records the generating program and its digest, the power-of-two exponent
$39997$, the bit length $39998$, the strict decimal power $12040$, the decimal
digit count $12041$, the largest partial quotient $973919$ at index $16442$,
and an arithmetic self-check of the same routine against the known expansions
of $e$ and $\pi$.  The certified statement is $q\ge2^{39997}$ and
$q>10^{12040}$.  Printing the bit length or the digit count as an exponent
overstates the bound by one power of two and one decimal order.

Each floor is finite, and extending either one raises the excluded range and
reproduces the same statement form.  Neither decides the problem.

\section{What cannot supply the missing input}\label{long251:sec:obstructions}

This section proves four boundaries.  Each says that a named class of
hypotheses about the prime gaps is compatible with a rational dyadic value, or
that a named event is too sparse for a named method.

\subsection{Bounded residue-preserving perturbations}

\begin{theorem}[bounded-perturbation obstruction]\label{long251:res:boundedperturbation}
Let $a_n$ be natural numbers with $\sum_{n\ge0}a_n2^{-(n+1)}$ convergent.  For
every integer $M\ge1$ and every cutoff $K$ there are digits
$\varepsilon_n\in\{0,1\}$, zero for $n<K$, such that
$\sum_{n\ge0}(a_n+M\varepsilon_n)2^{-(n+1)}$ is rational.
\end{theorem}

\begin{proof}
Write $A=\sum_{n\ge0}a_n2^{-(n+1)}$ and choose a rational
$r\in\bigl(A,A+M2^{-K}\bigr)$.  A binary expansion of $(r-A)/M$ has the form
$\sum_{n\ge K}\varepsilon_n2^{-(n+1)}$ with $\varepsilon_n\in\{0,1\}$; set
$\varepsilon_n=0$ for $n<K$.  The perturbed series converges and equals $r$.
\end{proof}

\noindent The
\lword{Erdos251/BoundedPerturbationCountermodel.lean}{193}{exists_bounded_perturbation}{construction}
is kernel checked.  The invariant-property consequence---a property shared by
every admissible perturbation cannot by itself force irrationality---is the
ordinary logical reading of that construction, rather than a separately named
Lean theorem.  The construction is also
checked at the actual consecutive prime gaps for every $M\ge1$ and every $K$,
the
\lword{Erdos251/BoundedPerturbationCountermodel.lean}{208}{exists_rational_bounded_perturbation_primeGap}{rational
perturbed prime-gap series}, with the perturbed gap lying in $[g_n,g_n+M]$ and
congruent to $g_n$ modulo $M$, the
\lword{Erdos251/BoundedPerturbationCountermodel.lean}{234}{perturbed_gap_bounds}{gap
bounds}.  Consequently no theorem about the size or the residue of prime gaps
that survives such a perturbation can suffice for Problem~\#251; that class
includes the bounded-gap and large-gap theorems cited in
Section~\ref{long251:sec:open} and equidistribution of $p_n$ modulo any fixed $q$, on
taking $M=2q$.  The perturbed digits need not remain prime gaps.

\subsection{Algebraic nonconcentration survives the rationalising perturbation}

The obvious reply to Theorem~\ref{long251:res:boundedperturbation} is to demand a
finer property of the actual gaps than size and residue.  A relevant such
property is Schlage-Puchta's Lemma~4: for every polynomial
$F\in\Z[x_0,\ldots,x_k]$ which does not vanish identically,
$F(g_n,\ldots,g_{n+k})\ne0$ for almost all $n$~\cite[Lemma~4,
pp.~5--6]{schlagepuchta2011}.  Say that an integer sequence $a$ has
\emph{fixed-block nonconcentration} when it satisfies that conclusion for
every $k\ge0$ and every nonzero $F\in\Z[x_0,\ldots,x_k]$.  The property
transfers across bounded perturbations with no independence or randomness
assumption.


\begin{proposition}[finite counting for shifted gap differences]\label{long251:res:shiftedcount}
Write $p_n$ for the primes indexed from $p_0=2$ and $g_n=p_{n+1}-p_n$.
For $h\ge2$ and $r\in\Z$, let $M_{h,r}(N)$ count $n<N$ with $g_{n+h}-g_n=r$.
Let $Q_{N,H,r}$ count triples $(x,d,s)$ with $x<p_N$,
$0<d<s\le H$, $d+r>0$, and all four integers
$x,x+d,x+s,x+s+d+r$ prime. Then, for every $N,H\ge0$,
\[
 (H+1)M_{h,r}(N)\le
 (h+1)p_{N+h+1}+(H+1)Q_{N,H,r}.
\]
\end{proposition}

\begin{proof}
Windows of total span at most $H$ inject into these four-prime
configurations; the sum of all length-$(h+1)$ spans is at most
$(h+1)p_{N+h+1}$, which bounds the number of larger spans.
\end{proof}

The \href{https://github.com/wcook04/plectis-erdos/blob/d058b9150218ae23d2cfd975580088d6167e9da8/ErdosProblems/Erdos251/ShiftedGapCountingR9.lean#L152}{finite shifted-count bound} is unconditional.
To deduce zero density one still needs, for every $\varepsilon>0$ and all
large $N$, a choice of $H$ making the right side less than
$\varepsilon N(H+1)$. The adjacent case $h=1$ instead requires the
corresponding three-prime estimate, because two of the four positions would
coincide. The all-shift implication records these analytic premises explicitly;
it supplies neither estimate and does not prove irrationality.

\begin{theorem}[nonconcentration is perturbation-stable]\label{long251:res:nonconcentration}
Let $a:\N\to\Z$ have fixed-block nonconcentration, let $E\subset\Z$ be finite,
and let $b_n=a_n+e_n$ with $e_n\in E$ for every $n$.  Then $b$ has fixed-block
nonconcentration.
\end{theorem}

\begin{proof}
Fix $k\ge0$ and a nonzero $F\in\Z[x_0,\ldots,x_k]$.  For a tuple
$\mathbf e=(e^{(0)},\ldots,e^{(k)})\in E^{k+1}$ put
$F_{\mathbf e}(x_0,\ldots,x_k)=F(x_0+e^{(0)},\ldots,x_k+e^{(k)})$.  The
substitution $x_i\mapsto x_i+e^{(i)}$ is a ring automorphism of
$\Z[x_0,\ldots,x_k]$ with inverse $x_i\mapsto x_i-e^{(i)}$, so
$F_{\mathbf e}\ne0$.  If $F(b_n,\ldots,b_{n+k})=0$ then
$F_{\mathbf e}(a_n,\ldots,a_{n+k})=0$ for the tuple
$\mathbf e=(e_n,\ldots,e_{n+k})$, whence
\[
 \{\,n:F(b_n,\ldots,b_{n+k})=0\,\}
 \;\subseteq\!\!
 \bigcup_{\mathbf e\in E^{k+1}}\!\!
 \{\,n:F_{\mathbf e}(a_n,\ldots,a_{n+k})=0\,\}.
\]
Each set on the right has density zero by hypothesis, and the union is over
the finite index set $E^{k+1}$, so a finite union of density-zero sets has
density zero.
\end{proof}

\begin{corollary}[nonconcentration does not force irrationality]\label{long251:res:nonconc-primes}
Fix $M\ge1$ and $K\ge0$, and let $b$ be the perturbed sequence supplied by
Theorem~\ref{long251:res:boundedperturbation} at the actual prime gaps.  Then
$\sum_{n\ge0}b_n2^{-(n+1)}$ is rational, $b_n=g_n$ for $n<K$,
$b_n-g_n\in\{0,M\}$ and $b_n\equiv g_n\pmod M$ for every $n$, $b$ has
fixed-block nonconcentration, and the cumulative sequence
$P_n=2+\sum_{i<n}b_i$ satisfies $p_n\le P_n\le p_n+Mn$ and hence
$P_n\sim n\log n$.
\end{corollary}

\begin{proof}
Everything except nonconcentration is Theorem~\ref{long251:res:boundedperturbation}
together with $\sum_{i<n}g_i=p_n-2$.  The perturbation takes values in the
fixed two-element set $E=\{0,M\}$, so Theorem~\ref{long251:res:nonconcentration}
applies to the actual gaps, which have fixed-block nonconcentration by
Schlage-Puchta's Lemma~4.
\end{proof}

Fixed-block polynomial nonconcentration, taken together with a prescribed
finite prefix of actual gaps, a pointwise bound $b_n\le g_n+M$, every residue
modulo $M$, positivity and cumulative growth at the prime-number-theorem scale,
is therefore compatible with a rational dyadic value.  Taking $M$ even also
preserves eventual evenness.  To preserve a prescribed modulus $q$ together
with parity, take $M=2q$, with the corresponding pointwise bound
$b_n\le g_n+2q$.  The boundary is exact: the terms $P_n$ are not asserted to be prime,
and the conclusion says nothing about constraints with a block length or a
polynomial family growing with the index.

\subsection{The producer is an event of density zero}

The same lemma bounds the event that the local certificate consumes.

\begin{theorem}[sparsity of the two-window event]\label{long251:res:sparse}
Fix $h\ge1$.  The set of $N\ge1$ at which the three hypotheses of
Theorem~\ref{long251:res:smallpair} hold for the actual prime gaps has density zero.
For the same $h$, the set of $N$ with $g_{N+h+1}=g_{N+1}$ also has density
zero.
\end{theorem}

\begin{proof}
For $N\ge1$ every gap involved is even, so $\delta_N=g_{N+h+1}-g_{N+1}$ is
even.  If $|\sigma_h(N)|<1$, $|\sigma_h(N+1)|<1$ and $\delta_N\ne0$, then
$\delta_N=2\sigma_h(N)-\sigma_h(N+1)$ lies in $(-3,3)$, so $\delta_N=\pm2$.
Apply Schlage-Puchta's Lemma~4 at index $n=N+1$ with $k=h$ to the two
polynomials $F_\pm(x_0,\ldots,x_h)=x_h-x_0\mp2$, neither of which vanishes
identically: each of $\{N:\delta_N=2\}$ and $\{N:\delta_N=-2\}$ has density
zero, and the event is contained in their union.  The second assertion is the
same lemma applied to $F(x_0,\ldots,x_h)=x_h-x_0$.
\end{proof}

The two halves point in opposite directions and both are useful.  The mismatch
hypothesis on its own holds for almost all $N$, which is strictly stronger than
the cofinal statement Proposition~\ref{long251:res:gap-nonperiodic} supplies.  The full
conjunction is confined to a set of density zero, because the two window
conditions force the mismatch to be exactly $\pm2$.  A producer for
Problem~\ref{long251:prob:smallpair} therefore has to be a cofinality statement about
a sparse set.  Consequently, this event cannot be supplied on a set of positive
lower density.  A sufficient producer would have to yield cofinally many
witnesses in a density-zero set; density zero alone does not exclude an
averaging argument capable of detecting such sparse witnesses.  The measurement reported
in Section~\ref{long251:sec:measurements} is consistent with this: the observed density
at $h=1$ declines like $1/\log p$, and the observed count up to $X$ grows like
$X/(\log X)^2$.

\subsection{Recurring gap values differing by two do not suffice}

Theorem~\ref{long251:res:signedwindow} reduces the producer to a $\pm2$ mismatch
together with a half-window condition.  It is tempting to hope that the
mismatch half is the substance, and that two even values differing by $2$, each
occurring infinitely often as a consecutive prime gap inside a common index
residue class, would suffice.  At the level of integer-digit recurrences that
implication is false, and it stays false under growth at the scale the primes
actually have.

\begin{theorem}[recurring values are not enough]\label{long251:res:polignacfail}
There is a sequence $(a_n)_{n\ge1}$ of positive even integers with the
following properties.  The values $2$ and $4$ each occur infinitely often at
indices divisible by every fixed $t\ge1$.  The sequence is unbounded, not
eventually periodic, and satisfies $a_n=O(\log n)$.  The series
$\sum_{n\ge1}a_n2^{-n}$ equals $6$, and every scaled tail
$\sum_{j\ge1}a_{N+j}2^{-j}$ is an integer, so every tail shift is integral.
The increasing odd sequence $P_n=3+\sum_{j\le n}a_j$ satisfies
$P_n\sim n\log n$.
\end{theorem}

\begin{proof}
Let logarithms be natural and put $b_n=2\lceil\tfrac12\log(n+64)\rceil$, so
that $b_n$ is even, $b_n\ge6$, and $b_n=\log(n+64)+O(1)$.  Since
$\tfrac12\log(n+65)-\tfrac12\log(n+63)<1$ for every $n\ge1$, consecutive
increments of $b$ lie in $\{0,2\}$ and $b_{n+1}-b_{n-1}\le2$.  Call an index
$n$ \emph{special} when $n=k!$ or $n=2\,k!$ for some $k\ge5$, and define
\[
 U_n=\begin{cases}
  2b_{n-1}-2,& n=k!,\ k\ge5,\\
  2b_{n-1}-4,& n=2\,k!,\ k\ge5,\\
  b_n,&\text{otherwise},
 \end{cases}
 \qquad a_n=2U_{n-1}-U_n\quad(n\ge1).
\]
The special indices are distinct and never adjacent, since $k!$ and $2\,k!$ are
even and differ by more than one from each other and from the neighbouring
special indices.  Each $U_n$ is an even integer by construction, so
each $a_n$ is an even integer.

At an ordinary index whose predecessor is also ordinary,
$a_n=2b_{n-1}-b_n=b_{n-1}-(b_n-b_{n-1})\ge6-2=4$.  At a special index the
predecessor is ordinary and $a_n=2b_{n-1}-(2b_{n-1}-r)=r$, which is $2$ at
$n=k!$ and $4$ at $n=2\,k!$.  Immediately after a special index carrying $r$,
\[
 a_{n+1}=2(2b_{n-1}-r)-b_{n+1}\ge4b_{n-1}-2r-(b_{n-1}+2)
 =3b_{n-1}-2r-2\ge8 .
\]
Every $a_n$ is therefore a positive even integer, and $U_n=O(\log n)$ gives
$a_n=O(\log n)$.  At ordinary indices with ordinary predecessors
$a_n\ge b_{n-1}-2\to\infty$, so $(a_n)$ is unbounded and hence not eventually
periodic.

The definition $a_n=2U_{n-1}-U_n$ is exactly $a_n2^{-n}=U_{n-1}2^{-(n-1)}-U_n2^{-n}$,
so for every $N\ge0$ and $m\ge1$
\[
 \sum_{j=1}^{m}\frac{a_{N+j}}{2^{\,j}}=U_N-\frac{U_{N+m}}{2^{\,m}} .
\]
Since $U_n=O(\log n)$ the endpoint tends to zero, so the tail at $N$ equals the
integer $U_N$; at $N=0$ the series equals $U_0=b_0=6$.  Every difference
$U_{N+h}-U_N$ is an integer, so every tail shift is integral.

For each $t$ and every sufficiently large $k$ we have $t\mid k!$, hence
$t\mid2\,k!$, and $a_{k!}=2$ while $a_{2k!}=4$: both values recur infinitely
often at indices congruent to $0$ modulo $t$.

Finally, summing $a_j=2U_{j-1}-U_j$ gives
$\sum_{j=1}^{n}a_j=2U_0+\sum_{j=1}^{n-1}U_j-U_n$.  The baseline
$\sum_{j<n}b_j=n\log n+O(n)$.  The special indices up to $n$ number
$O(\log n/\log\log n)$ and each changes the summand by $O(\log n)$, so their
total contribution is $O((\log n)^2)=o(n)$.  Hence
$P_n=3+\sum_{j\le n}a_j\sim n\log n$, and $P_n$ is odd and increasing.
\end{proof}

The factorial schedule above proves recurrence in the zero residue class.
The following separate construction strengthens this to every residue class.


\subsection{A complete countermodel in every residue class}
\label{long251:res:all-residue-log-countermodel}
There is a synthetic integer sequence $(a_n)_{n\ge1}$ with all of the
following properties.  Each $a_n$ is positive and divisible by $2$, and
\[
 a_n\le 4\log(n+1)+24.
\]
For every $t\ge1$, every $0\le r<t$, and every cutoff $N$, there are
$i,j\ge N$ such that
\[
 i\equiv j\equiv r\pmod t,\qquad a_i=2,\quad a_j=4.
\]
For every integer $B$ and cutoff $N$, some $n\ge N$ satisfies $a_n>B$.
For every $h\ge1$ there is no cutoff after which $a_{n+h}=a_n$ for all $n$.
Nevertheless, all the complete tails
\[
 U_N=\sum_{j\ge1}\frac{a_{N+j}}{2^j}\quad(N\ge0)
\]
are integers, $U_0=6$, and $U_{N+h}-U_N\in\mathbb Z$ for every $N,h\ge0$.
The sequence
\[
 P_N=3+\sum_{j=1}^{N}a_j
\]
is strictly increasing and odd, with $P_N/(N\log N)\to1$.
These are synthetic positions; no $P_N$ is asserted to be prime.

To obtain every residue class, enumerate triples consisting of a positive
modulus, a residue, and a repetition index.  Choose the corresponding
centres $c_k$ in the prescribed classes with $c_0\ge100$,
$c_{k+1}\ge c_k+3$ and $c_k\ge2^{k^2}$ for $k\ge1$.
The parity of the repetition index selects the desired value $v_k=2$ or $4$.
Use the even baseline
\[
 b_n=6+2\left\lfloor\frac{\log(n+1)}2\right\rfloor,
 \qquad
 U_n=\begin{cases}2b_{n-1}-v_k,&n=c_k,\\b_n,&n\notin\{c_k:k\ge0\},\end{cases}
\]
and define $a_n=2U_{n-1}-U_n$.  The separation of the centres and the
baseline increment bound give positivity and the displayed logarithmic
bound.  Finite telescoping leaves $U_{N+m}/2^m$, which tends to zero;
thus these carries are the complete tails.  The sparse changes to the
baseline have negligible contribution after division by $N\log N$,
which gives the stated growth of $P_N$.

The full statement is \href{https://github.com/wcook04/plectis-erdos/blob/d4fed71423840f70f10edf27b9ad27c22fc4f49a/ErdosProblems/Erdos251/AllResidueLogarithmicR9.lean#L503}{kernel-checked in Lean}.


This rules out an inference from these coefficient properties alone to
nonintegral tail shifts.  It gives no counterexample to the prime-gap problem
and does not reproduce the extreme large gaps of the primes.

The consequence is exact.  Positivity, evenness, unboundedness, non-eventual
periodicity, a pointwise logarithmic bound, cumulative growth at the
prime-number-theorem scale, and the recurrence of two even values differing by
$2$ inside every index residue class are jointly compatible with an integral
tail at every index.  Any route through that input must use a property of the
consecutive primes beyond this list.  The sequence $(a_n)$ is synthetic and the
$P_n$ are not asserted to be prime.  The pointwise logarithmic bound also
precludes the known large-gap behaviour, so the theorem does not speak to
hypotheses that use extreme gaps.

\subsection{Two further boundaries}

\begin{proposition}[quadratic polynomial-shift countermodel]\label{long251:res:polynomialcountermodel}
Put $c_n=2(n^2+4n+2)$ and $U_n=2(n+4)^2$.  Then $c_n$ is positive, even and
strictly increasing, $U_{n+1}=2U_n-c_{n+1}$, every shift $U_{N+h}-U_N$ is
integral, $c_{n+1}-c_n=4n+10$ is never $\pm2$, and
\[
 \sum_{j\ge1}\frac{c_j}{2^{\,j}}=32 .
\]
\end{proposition}

\begin{proof}
Direct expansion gives the recurrence, and the finite telescope is
$\sum_{j=1}^{n}c_j2^{-j}=32-2(n+4)^22^{-n}$, whose last term tends to zero.
The remaining assertions follow from the integer values of $U_n$ and from
$c_{n+1}-c_n=4n+10\ge10$.
\end{proof}

\noindent The series is the one starting at $j=1$, whose initial tail state is
$U_0=32$; under the opening convention $\sum_{n\ge0}c_n2^{-(n+1)}$ the same
word has value $18$, and the two ranges must be kept apart.  The formal
construction checks the
\lword{Erdos251/PrimeGapDyadicTail.lean}{1596}{polynomialTailOrbit_recurrence}{recurrence},
\lword{Erdos251/PrimeGapDyadicTail.lean}{1632}{polynomialGapWord_strictMono}{strict
growth},
\lword{Erdos251/PrimeGapDyadicTail.lean}{1654}{polynomialTailOrbit_shift_integral}{integrality
of every shift},
\lword{Erdos251/PrimeGapDyadicTail.lean}{1644}{polynomialGapWord_no_adjacent_two_digit}{exclusion
of adjacent differences of size two}, and the
\lword{Erdos251/PolynomialGapSeriesValue.lean}{92}{tsum_polynomialGapDyadicTerm_eq}{value
$32$ for the shifted series}, whose summand is $c_{n+1}/2^{\,n+1}$; the
rational value is recorded at
\lword{Erdos251/PolynomialGapSeriesValue.lean}{109}{not_irrational_tsum_polynomialGapDyadicTerm}{line~109}.
Positivity, parity, strict growth, unboundedness and nonperiodicity are
therefore jointly compatible with rationality, and the adjacent trigger of
Theorem~\ref{long251:res:signedwindow} fails at every index.  The telescoping mechanism
is the one Kova\v{c} used against the variable-denominator conjecture printed
under the problem~\cite[Theorem~1 and proof, pp.~1--2]{kovac2026}; the word
carries no priority claim and is not a result about Problem~\#251.

Rationality also fails to force the coefficients to repeat.  Let $K:\N\to\Q$ be
arbitrary and put $\kappa_n=2K_n-K_{n+1}$: the
\lword{Erdos251/PrimeGapDyadicTail.lean}{1129}{carryCoeff}{carry coefficient}.

\begin{proposition}[exact telescoping]\label{long251:res:telescope}
For every $n\ge0$,
$\sum_{i=0}^{n-1}\kappa_i2^{-(i+1)}=K_0-K_n2^{-n}$.
\end{proposition}

\begin{proof}
A routine induction; the added term is
$(2K_n-K_{n+1})2^{-(n+1)}=K_n2^{-n}-K_{n+1}2^{-(n+1)}$.
\end{proof}

\noindent Formalised as the
\lword{Erdos251/PrimeGapDyadicTail.lean}{1137}{carryPartialSum_eq}{carry
telescoping identity}.  Taking $K_0=\tfrac52$ and $K_n=2n+2$ for $n\ge1$ gives
$\kappa_0=1$ and $\kappa_n=2n$, so the coefficients $1,2,4,6,8,\ldots$ are
positive, even after the first term, unbounded and not eventually periodic,
while their dyadic sum is $\tfrac52$.  Rationality alone therefore cannot imply
eventual periodicity even for a positive, parity-correct integer coefficient
sequence, and no argument may play periodicity of a rational value against
Proposition~\ref{long251:res:gap-nonperiodic}.

Finally, two escape tests that look like independent producers are equivalent
to the conclusion they were introduced to supply.  For fixed $h$ write
$\delta_N=g_{N+h+1}-g_{N+1}$, $D_N=\sigma_h(N)$ and
$B_{h,N,r}=\sum_{i=0}^{r-1}2^{\,r-1-i}\delta_{N+i}$, so that
$D_{N+r}=2^{r}D_N-B_{h,N,r}$, and call the data-dependent affine condition
$\mathcal A_{N,r}$ the assertion $D_{N+r}\in-B_{h,N,r}+2^{\,r+1}\Z$.

\begin{theorem}[affine and fixed-lattice circularity]\label{long251:res:affinecollapse}
For every rational dyadic tail recurrence and all $h,N,r\ge0$,
$\mathcal A_{N,r}$ holds if and only if $D_N\in2\Z$.  Consequently, if every
$\delta_N$ is even, cofinal failure of $\mathcal A$ is equivalent to
nonintegrality of $D_N$ for arbitrarily large $N$.  The same equivalence holds
for the fixed-lattice replacement, under any bound $|D_N|\le b(N)$ satisfying
$2b(N+r)q<2^{r}$ for some $r$ at every $N$ and every positive integer $q$.
\end{theorem}

\begin{proof}
Substituting $D_{N+r}=2^{r}D_N-B_{h,N,r}$ into $\mathcal A_{N,r}$ cancels the
observed block and leaves $D_N=2z$.  After one recurrence step, evenness of
$\delta_N$ identifies even integrality at $N+1$ with ordinary integrality at
$N$; the reverse implication chooses a nonintegral $D_N$ and depth $r=0$.  For
the fixed-lattice form, eventual integrality and the block identity give
$B_{h,N,r}-2^{r}z=-D_{N+r}$, so a strict separation exceeding $b(N+r)$ is
contradictory; conversely a nonintegral $D_N$ with reduced denominator $q$ has
distance at least $1/q$ from every integer, and scaling that separation by
$2^{r}$ at a legal depth gives the required strict inequality.
\end{proof}

\noindent These are the
\lword{Erdos251/AffineCylinderCollapse.lean}{91}{ratAffinePowTwo_iff_evenIntegral}{pointwise
affine identity}, its
\lword{Erdos251/AffineCylinderCollapse.lean}{162}{cofinal_affinePowTwo_escape_iff_not_eventuallyIntegral}{cofinal
form}, and the
\lword{Erdos251/AffineCylinderCollapse.lean}{308}{cofinal_blockResidue_escape_iff_not_eventuallyIntegral}{fixed-lattice
equivalence}.  The apparent affine hierarchy contains no depth-dependent
information, and the fixed lattice removes the data dependence at the cost of a
growth condition under which rational denominator separation already forces
every required escape.  Neither is an independent source of information about
consecutive primes.

\section{What is measured}\label{long251:sec:measurements}

Two finite computations report on the two producers.  Each covers one bounded
range and supplies no cofinal statement.

Over the $6\,841\,648$ primes below $1.2\times10^{8}$ and offsets
$h=1,\ldots,16$, the reduced two-window event of
Theorem~\ref{long251:res:signedwindow} occurs at every offset, with density between
$0.00418$ and $0.008248$ and the two signs near balanced; at $h=1$ there are
$56\,427$ occurrences among $6\,841\,564$ candidate indices, the last at the
prime $119\,995\,753$.  Rescaled by $\log p$ the band densities run from
$0.17671$ down to $0.13148$ with decelerating drift, so the count up to $X$
still grows like $X/(\log X)^2$, which is the rate
Theorem~\ref{long251:res:sparse} permits.  The tails in this scan are double-precision
floating-point, so an individual hit is a numerical observation and carries no
certificate; the median distance from the shift to the nearer window boundary
is $0.127$ and the worst case over $34\,000$ hits is $6.1\times10^{-7}$, eight
orders above the working resolution, and $3\,200$ sampled hits across
$h=1,\ldots,8$ were rechecked against the literal three-part statement with no
violations.  Proposition~\ref{long251:res:finite-smallpair} is the one pair certified
by exact integer arithmetic.  The receipt is
\texttt{state/formal\_math/probes/erdos251\_adjacent\_mismatch\_receipt.json}.

Over the $1\,270\,607$ primes below $2\times10^{7}$, and for every modulus
$t\le20$ and every residue class modulo $t$, the free-pair event of
Theorem~\ref{long251:res:freepair} has witnesses running to the last usable index: the
leanest class at $t=20$ carries $90\,150$ witnesses and the latest witness sits
at index $1\,270\,519$ of $1\,270\,540$.  Every recorded witness satisfies the
stronger two-window event, so it certifies both producers at once.  What the
scan does not supply is uniformity in $t$ and in the cutoff.  The receipt is
\texttt{state/formal\_math/erdos257\_period\_noncollapse/}\allowbreak
\texttt{erdos251\_free\_pair\_receipt.json}.

\section{The remaining obligation}\label{long251:sec:open}

Problem~\#251 is open.  The exact unresolved condition, equivalent to
irrationality by Theorem~\ref{long251:res:escape-irrational} and the escape
equivalence for the actual tail, is the following.

\begin{problem}[universal prime-gap shift escape]\label{long251:prob:escape}
For every $h\ge1$ and every $N_0$, prove that some $N\ge N_0$ satisfies
\[
 \sum_{j\ge1}\frac{g_{N+h+j}-g_{N+j}}{2^{\,j}}\notin\Z .
\tag{8.1}\label{long251:eq:shift-escape}
\]
\end{problem}

\begin{problem}[cofinal adjacent small mismatch]\label{long251:prob:smallpair}
For every $h\ge1$ and every $N_0$, prove that some $N\ge N_0$ satisfies
\[
 \bigl|\sigma_h(N)\bigr|<1,\qquad
 \bigl|\sigma_h(N+1)\bigr|<1,\qquad
 g_{N+h+1}\ne g_{N+1} .
\tag{8.2}\label{long251:eq:smallpair}
\]
\end{problem}

Corollary~\ref{long251:res:smallpair-real} makes Problem~\ref{long251:prob:smallpair}
sufficient for Problem~\ref{long251:prob:escape}, and Theorem~\ref{long251:res:freepair}
supplies an equivalent target in which the offset is free.  The three
formulations are equivalent or ordered by implication;
none of them is an analytic saving on its own.

What the proved obstructions leave.  By Theorem~\ref{long251:res:boundedperturbation}
no hypothesis stable under bounded residue-preserving perturbation can work;
by Corollary~\ref{long251:res:nonconc-primes} adding fixed-block polynomial
nonconcentration to that list does not repair it; the construction in
Section~\ref{long251:res:all-residue-log-countermodel} shows that recurring gap
values differing by $2$ in every residue class do not work at the level of
integer recurrences; and by
Theorem~\ref{long251:res:sparse} the event in \eqref{long251:eq:smallpair} has density zero, so
it must be produced cofinally on a sparse set.  Isolated small gaps, isolated
large gaps and average gap estimates are insufficient.  Any prescribed finite
prefix can be preserved while the complete sum is rationalised, so that prefix
alone cannot force irrationality or a cofinal small-tail producer.  A finite
block together with a proved tail bound can still certify one window: both
inequalities in \eqref{long251:eq:smallpair} contain the complete infinite
continuation, and Proposition~\ref{long251:res:explicit-remainder} supplies the
corresponding finite reduction.  Nor does
parity help, since after the first gap every $g_n$ is even.

The finite reduction is available.  Proposition~\ref{long251:res:explicit-remainder}
turns each window condition into an integer comparison at truncation length
$L$, and for fixed $h$ and $\varepsilon>0$ the choice
$L=\lceil(4+\varepsilon)\log_2(N+2)\rceil$ makes $E_{h,N,L}=O_h(N^{-\varepsilon})$.
A prescribed positive margin $\eta_h$ therefore converts \eqref{long251:eq:smallpair}
into a statement about the joint distribution modulo powers of two of the
finite block $(g_{N+h+1}-g_{N+1},\ldots,g_{N+h+L}-g_{N+L})$ over a window of
logarithmic length.  Deciding one instance costs $O(\log N)$ gaps together with
the elementary prime bound; producing them cofinally, with uniformity in $h$
and in the cutoff, is the open work.

The strongest published results on prime gaps address a different shape of
question.  Zhang's bounded-gap theorem~\cite[Theorem~1, p.~1122]{zhang2014}
produces infinitely many bounded consecutive-prime gaps.  Maynard bounds
$\liminf_n(p_{n+m}-p_n)$ for every fixed $m$, giving bounded clusters of every
fixed size~\cite[Theorem~1.1, p.~384]{maynard2015}, with the explicit
unconditional bound $\liminf_n(p_{n+1}-p_n)\le600$~\cite[Theorem~1.3,
p.~385]{maynard2015}.  The large-gap theorem of Ford, Green, Konyagin, Maynard
and Tao bounds the largest single consecutive-prime gap below
$X$~\cite[Theorem~1, p.~66]{fgkmt2018}.  A theorem giving bounded clusters at
each fixed cluster size does not by itself supply a weighted condition on
windows whose length grows with the basepoint, and none of these results is
used as a proof input here.

Conditionally the picture is different.  Tao's comment names uniform
quantitative prime-tuples control of about $\log\log n$ consecutive gaps as the
plausible route~\cite[comment of 7 October 2025]{erdosproblems251thread}, and
Land's draft carries that out under Kuperberg's uniform prime-tuples
conjecture~\cite{land2026}.  The obstructions of
Section~\ref{long251:sec:obstructions} say which unconditional substitutes cannot
replace that hypothesis.

\paragraph{What remains to be formalised.}
The two finite certificates of
Propositions~\ref{long251:res:explicit-remainder} and~\ref{long251:res:finite-smallpair}, the
continued-fraction exclusion of Theorem~\ref{long251:res:cfexclusion}, the
nonconcentration transfer of Theorem~\ref{long251:res:nonconcentration} with its
corollary, the sparsity statement of Theorem~\ref{long251:res:sparse}, and the
construction of Theorem~\ref{long251:res:polignacfail} are ordinary mathematical proofs
with exact arithmetic replay and no Lean declaration.  Everything else linked
above is checked by the kernel at the pinned commit.

\pdfbookmark[1]{Statements and declarations}{statements-declarations}
\subsection*{Statements and declarations}

\paragraph{Artefact and data availability.} The
\href{\sourceurl}{pinned formal-source revision} contains the Lean sources, the
fixed toolchain, and the library manifest used in the verification.  Proof
authority rests in those pinned sources.  The present text is exposition and
navigation.

Lean checks each proof term against the fixed library version, and the sources
linked here contain no proof placeholders and no project-defined axioms; Lean
does not authorise the exposition, the citation choices, or the
interpretation, for which the author remains responsible.  The two receipts
named in Section~\ref{long251:sec:measurements} and the continued-fraction receipt of
Theorem~\ref{long251:res:cfexclusion} record their generating programs with digests
and are held in the private research development.

\paragraph{Funding and competing interests.} This work received no external
funding.  The author declares no competing interests.

\paragraph{Acknowledgements.} The problem numbering and status follow the
Erd\H{o}s Problems catalogue maintained by Thomas Bloom~\cite{erdosproblems}.
The logarithmic-scale construction of Theorem~\ref{long251:res:polignacfail} and the
transfer argument of Theorem~\ref{long251:res:nonconcentration} arose in an external
review of an earlier draft of this note; both proofs printed here were checked
independently.

\appendix
\section{An elementary polynomial bound for the primes}\label{long251:app:prime-bound}

We prove $p_n\le1250(n+1)^4$ for every $n\ge0$.  Let $\pi(x)$ count the primes
at most $x$.  For an integer $m\ge4$,
\begin{equation}\label{long251:eq:binomial-bound}
 4^m<m\binom{2m}{m}\le m(2m)^{\pi(2m)} .
\end{equation}
For the first inequality, $m\binom{2m}{m}/4^{m}$ equals $70/64$ at $m=4$ and
its ratio at successive indices is $(2m+1)/(2m)>1$.  For the second, the
exponent of a prime $\ell$ in $\binom{2m}{m}$ is
$\sum_{k\ge1}\bigl(\lfloor2m/\ell^{k}\rfloor-2\lfloor m/\ell^{k}\rfloor\bigr)$,
each summand is $0$ or $1$, and every summand with $\ell^{k}>2m$ vanishes;
hence each prime power dividing $\binom{2m}{m}$ is at most $2m$, and there are
$\pi(2m)$ of them.

Now fix $n\ge0$, put $x=n+5$ and $m=x^4\ge625$, and suppose $\pi(2m)\le n$.
Since $x\le2^{x}$ and $n+4(n+1)x=4x^2-15x-5\le2x^4$,
\[
 m(2m)^{n}=2^{n}x^{4(n+1)}
 \le2^{\,n+4(n+1)x}\le2^{\,2x^4}=4^{m},
\]
contradicting \eqref{long251:eq:binomial-bound}.  Hence $\pi(2m)>n$, so at least $n+1$
primes lie below $2m$ and therefore
\[
 p_n\le2(n+5)^4\le1250(n+1)^4 ,
\]
the last step because $(n+5)\le5(n+1)$ for $n\ge0$.  The same bound is checked
by the kernel at
\lword{Erdos251/PrimeGapDyadicTail.lean}{360}{prime0_le_polynomial}{line~360}.

\section{Integer certificates}\label{long251:app:certificates}

The following calculation uses trial division and integer arithmetic only.  It
reproduces the two rows of Proposition~\ref{long251:res:finite-smallpair} and every
inequality used in Theorem~\ref{long251:res:denominatorfloor}.  The four literals are
the certificate the Lean kernel decides in
\texttt{KernelDenominatorFloor.lean}; their use here requires no
floating-point approximation.  Adjacent quoted strings are concatenated by
Python.
{\footnotesize
\begin{verbatim}
from math import isqrt

primes = [n for n in range(2, 10000)
          if all(n % d for d in range(2, isqrt(n) + 1))]
gaps = [q - p for p, q in zip(primes, primes[1:])]
P = lambda x: x**4 + 8*x**3 + 36*x*x + 104*x + 150
Q = 2**40
expected = [(-662838684750, 11764181250),
            (873345886050, 12805761250)]
for N, target in zip((2, 3), expected):
    D = sum((gaps[N+1+j] - gaps[N+j])*2**(40-j)
            for j in range(1, 41))
    B = 1250*(P(N+43) + P(N+42))
    distance = min(D % Q, Q - D % Q)
    if (D, B) != target or not (abs(D)+B < Q and B < distance):
        raise ArithmeticError("tail certificate failed")
if gaps[4] == gaps[3]:
    raise ArithmeticError("gap mismatch failed")

u = int(
    "8065641857152652932176019632186898003271162829171466334827"
    "3083607794415278717445033509407855988903369988525550746159"
    "7355889792250084202344821020139160956663658789718168152662"
    "0217"
)

v = int(
    "2194945124413663232143970924541263312422069524635615360518"
    "4247077351958221816830720189289904831662955084392690248683"
    "1216291723988537733235173040607254496838530213867781442335"
    "1745"
)

up = int(
    "6539437101498162626882411892470905222108260008568555445975"
    "3026163315521784668689909712759876562484620059098138417469"
    "5232839888185316374272333277389611483117334000493867584923"
    "912"
)

vp = int(
    "1779611075789866551198427241621709630120136323281598176637"
    "8490605249229735501968764904036152970729491656650873938580"
    "6203025466313389810291243786005499164537499383612081805160"
    "767"
)

c = len(primes)
A = sum(p*2**(c-1-i) for i, p in enumerate(primes))
checks = (c == 1229,
          u > 0 and v > 0 and up > 0 and vp > 0,
          up*v - u*vp == 1,
          u*2**c < A*v,
          (2*A + 5000*(c+1)**4)*vp < up*2**(c+1),
          v+vp >= 2**589,
          2**589 > 10**177)
if not all(checks):
    raise ArithmeticError("denominator certificate failed")
print("Both tail rows and the denominator floor verified.")
\end{verbatim}
}

\section{Guide to the formal sources}\label{long251:app:index}

Each linked phrase opens its Lean declaration at the pinned source
revision~\commitshort.  The summation-by-parts and prime-bound declarations
are prime-specific.  Section~\ref{long251:sec:tail} is stated for arbitrary integer
digits and arbitrary rational or real orbits, while
Section~\ref{long251:sec:local-certificate} records the actual-gap specialisation.  The
concrete prime-gap tail, its unconditional convergence, its recurrence and the
real-to-rational scaled-tail bridge are defined and checked in
\texttt{RealPrimeGapTail.lean} and \texttt{PrimeGapDyadicTail.lean}; the
free-pair lattice and its actual-series equivalence in
\texttt{FreePairReduction.lean}; the kernel-decided denominator certificate in
\texttt{KernelDenominatorFloor.lean}; the rational countermodels in
\texttt{BoundedPerturbationCountermodel.lean} and
\texttt{PolynomialGapSeriesValue.lean}; and the two circularity theorems in
\texttt{AffineCylinderCollapse.lean}, with the signed normal form in
\texttt{AffineShiftEscape.lean} and the lcm diagonal in
\texttt{OrderLatticeDiagonal.lean}.

% The exhaustive source manifest is retained in the authored TeX for automated
% validation, and kept outside the printed note.
\iffalse

\begin{tabular}{@{}ll@{}}
informal statement & linked Lean source\\
zero-based primes & \lrefx{Erdos251/PrimeGapDyadicTail.lean}{40}{prime0}\\
zero-based gaps & \lrefx{Erdos251/PrimeGapDyadicTail.lean}{44}{primeGap0}\\
first gap value & \lrefx{Erdos251/PrimeGapDyadicTail.lean}{47}{primeGap0_zero}\\
second gap value & \lrefx{Erdos251/PrimeGapDyadicTail.lean}{51}{primeGap0_one}\\
unbounded actual prime gaps & \lrefx{Erdos251/PrimeGapDyadicTail.lean}{57}{exists_primeGap0_gt}\\
dyadic partial sum & \lrefx{Erdos251/PrimeGapDyadicTail.lean}{101}{dyadicPartialSumQ}\\
displayed normalisation & \lrefx{Erdos251/PrimeGapDyadicTail.lean}{106}{prime0DisplayedPartialSumQ}\\
factor-of-two identity & \lrefx{Erdos251/PrimeGapDyadicTail.lean}{111}{prime0DisplayedPartialSumQ_eq_two_mul}\\
dyadic difference sum & \lrefx{Erdos251/PrimeGapDyadicTail.lean}{121}{dyadicDifferencePartialSumQ}\\
partial-sum step & \lrefx{Erdos251/PrimeGapDyadicTail.lean}{124}{dyadicPartialSumQ_succ}\\
difference-sum step & \lrefx{Erdos251/PrimeGapDyadicTail.lean}{129}{dyadicDifferencePartialSumQ_succ}\\
summation by parts & \lrefx{Erdos251/PrimeGapDyadicTail.lean}{138}{dyadicPartialSumQ_eq_start_add_differences}\\
monotonicity step & \lrefx{Erdos251/PrimeGapDyadicTail.lean}{152}{prime0_mono_step}\\
gap cast identity & \lrefx{Erdos251/PrimeGapDyadicTail.lean}{156}{primeGap0_cast}\\
gap partial sum & \lrefx{Erdos251/PrimeGapDyadicTail.lean}{161}{primeGapPartialSumQ}\\
prime-gap summation by parts & \lrefx{Erdos251/PrimeGapDyadicTail.lean}{172}{prime0_dyadic_summation_by_parts}\\
prime dyadic term & \lrefx{Erdos251/PrimeGapDyadicTail.lean}{183}{primeDyadicTerm}\\
displayed prime dyadic term & \lrefx{Erdos251/PrimeGapDyadicTail.lean}{188}{primeDisplayedDyadicTerm}\\
prime-gap dyadic term & \lrefx{Erdos251/PrimeGapDyadicTail.lean}{192}{primeGapDyadicTerm}\\
infinite factor-of-two identity & \lrefx{Erdos251/PrimeGapDyadicTail.lean}{196}{primeDisplayedDyadicTerm_eq_two_mul}\\
gap discrete derivative & \lrefx{Erdos251/PrimeGapDyadicTail.lean}{202}{primeGapDyadicTerm_eq}\\
polynomial prime bound & \lrefx{Erdos251/PrimeGapDyadicTail.lean}{360}{prime0_le_polynomial}\\
prime-series summability & \lrefx{Erdos251/PrimeGapDyadicTail.lean}{379}{summable_primeDyadicTerm}\\
gap summability transfer & \lrefx{Erdos251/PrimeGapDyadicTail.lean}{385}{summable_primeGapDyadicTerm_of_summable_primeDyadicTerm}\\
infinite gap identity & \lrefx{Erdos251/PrimeGapDyadicTail.lean}{404}{tsum_primeDyadicTerm_eq_two_add_primeGap}\\
unconditional gap identity & \lrefx{Erdos251/PrimeGapDyadicTail.lean}{427}{tsum_primeDyadicTerm_eq_two_add_primeGap_unconditional}\\
normalised irrationality equivalence & \lrefx{Erdos251/PrimeGapDyadicTail.lean}{435}{irrational_tsum_primeDyadicTerm_iff_primeGap}\\
displayed infinite identity & \lrefx{Erdos251/PrimeGapDyadicTail.lean}{444}{tsum_primeDisplayedDyadicTerm_eq_four_add_two_primeGap}\\
displayed irrationality equivalence & \lrefx{Erdos251/PrimeGapDyadicTail.lean}{459}{irrational_tsum_primeDisplayedDyadicTerm_iff_primeGap}\\
tail recurrence & \lrefx{Erdos251/PrimeGapDyadicTail.lean}{482}{DyadicTailRecurrence}\\
rational actual-tail state & \lrefx{Erdos251/PrimeGapDyadicTail.lean}{489}{rationalPrimeGapTailState}\\
hypothetical-rational tail representation & \lrefx{Erdos251/PrimeGapDyadicTail.lean}{510}{exists_rationalPrimeGapTailState_representation_of_not_irrational}\\
rational actual-tail recurrence & \lrefx{Erdos251/PrimeGapDyadicTail.lean}{527}{rationalPrimeGapTailState_recurrence}\\
shift & \lrefx{Erdos251/PrimeGapDyadicTail.lean}{544}{tailShift}\\
shift step identity & \lrefx{Erdos251/PrimeGapDyadicTail.lean}{563}{tailShift_succ}\\
integrality & \lrefx{Erdos251/PrimeGapDyadicTail.lean}{575}{RatIntegral}\\
tail block & \lrefx{Erdos251/PrimeGapDyadicTail.lean}{647}{dyadicTailBlock}\\
iterated block identity & \lrefx{Erdos251/PrimeGapDyadicTail.lean}{653}{tail_iterate_eq_pow_mul_sub_block}\\
scaled shift identity & \lrefx{Erdos251/PrimeGapDyadicTail.lean}{667}{tailShift_eq_scaled_sub_block}\\
integer-shift invariance & \lrefx{Erdos251/PrimeGapDyadicTail.lean}{677}{ratIntegral_sub_int_iff}\\
Euler multiplier & \lrefx{Erdos251/PrimeGapDyadicTail.lean}{695}{ratIntegral_totientMultiplier_of_odd_den}\\
power-of-two/odd denominator scaling & \lrefx{Erdos251/PrimeGapDyadicTail.lean}{723}{ratIntegral_scaled_of_den_eq_pow_two_mul_odd}\\
multiplier at the two-adic exponent & \lrefx{Erdos251/PrimeGapDyadicTail.lean}{756}{ratIntegral_multiplier_tail_at_twoExponent}\\
integral-shift criterion & \lrefx{Erdos251/PrimeGapDyadicTail.lean}{802}{tailShift_integral_iff_scaledTail}\\
fixed-denominator eventual shift & \lrefx{Erdos251/PrimeGapDyadicTail.lean}{830}{rationalPrimeGapTailShift_eventuallyIntegral_of_den_eq}\\
eventual integral shift & \lrefx{Erdos251/PrimeGapDyadicTail.lean}{853}{rationalPrimeGapTailShift_eventuallyIntegral}\\
totient shift & \lrefx{Erdos251/PrimeGapDyadicTail.lean}{877}{tailShift_integral_totient_of_odd_den}\\
one-step propagation & \lrefx{Erdos251/PrimeGapDyadicTail.lean}{888}{tailShift_integral_succ}\\
propagation to every later index & \lrefx{Erdos251/PrimeGapDyadicTail.lean}{900}{tailShift_integral_add}\\
adjacent small-shift obstruction & \lrefx{Erdos251/PrimeGapDyadicTail.lean}{979}{tailShift_not_both_integral_of_small_pair_of_digit_ne}\\
cofinal small-mismatch theorem & \lrefx{Erdos251/PrimeGapDyadicTail.lean}{1006}{not_eventuallyIntegralTailShift_of_cofinal_small_mismatch}\\
prime gaps not eventually periodic & \lrefx{Erdos251/PrimeGapDyadicTail.lean}{1023}{primeGap0_not_eventually_periodic}\\
shift not eventually small & \lrefx{Erdos251/PrimeGapDyadicTail.lean}{1097}{rationalPrimeGapTail_has_positive_shift_not_eventually_small}\\
actual-prime-gap specialisation & \lrefx{Erdos251/PrimeGapDyadicTail.lean}{1112}{primeGapTailShift_not_eventuallyIntegral_of_cofinal_small_mismatch}\\
carry coefficient & \lrefx{Erdos251/PrimeGapDyadicTail.lean}{1129}{carryCoeff}\\
carry partial sum & \lrefx{Erdos251/PrimeGapDyadicTail.lean}{1133}{carryPartialSum}\\
carry telescoping identity & \lrefx{Erdos251/PrimeGapDyadicTail.lean}{1137}{carryPartialSum_eq}\\
natural carry coefficient & \lrefx{Erdos251/PrimeGapDyadicTail.lean}{1177}{natCarryCoeff}\\
natural-carry cast & \lrefx{Erdos251/PrimeGapDyadicTail.lean}{1180}{natCarryCoeff_cast}\\
exact denominator recurrence & \lrefx{Erdos251/PrimeGapDyadicTail.lean}{1217}{tail_den_succ}\\
odd-denominator persistence & \lrefx{Erdos251/PrimeGapDyadicTail.lean}{1226}{tail_den_succ_eq_of_odd}\\
even-denominator halving & \lrefx{Erdos251/PrimeGapDyadicTail.lean}{1236}{tail_den_succ_eq_half_of_even}\\
integral shifts by denominator & \lrefx{Erdos251/PrimeGapDyadicTail.lean}{1279}{tailShift_integral_iff_den_dvd_mersenne}\\
power-of-two congruence & \lrefx{Erdos251/PrimeGapDyadicTail.lean}{1291}{tailShift_integral_iff_two_pow_modEq_one}\\
one-shift rationality classifier & \lrefx{Erdos251/PrimeGapDyadicTail.lean}{1479}{not_irrational_initial_iff_exists_integral_positive_tailShift}\\
eventual-shift rationality classifier & \lrefx{Erdos251/PrimeGapDyadicTail.lean}{1525}{not_irrational_initial_iff_exists_eventually_integral_positive_tailShift}\\
pointwise irrationality classifier & \lrefx{Erdos251/PrimeGapDyadicTail.lean}{1551}{irrational_initial_iff_all_positive_tailShifts_nonintegral}\\
cofinal irrationality classifier & \lrefx{Erdos251/PrimeGapDyadicTail.lean}{1572}{irrational_initial_iff_cofinalNonintegralTailShifts}\\
polynomial countermodel recurrence & \lrefx{Erdos251/PrimeGapDyadicTail.lean}{1596}{polynomialTailOrbit_recurrence}\\
polynomial countermodel strict growth & \lrefx{Erdos251/PrimeGapDyadicTail.lean}{1632}{polynomialGapWord_strictMono}\\
polynomial countermodel excludes the adjacent $\pm2$ trigger & \lrefx{Erdos251/PrimeGapDyadicTail.lean}{1644}{polynomialGapWord_no_adjacent_two_digit}\\
all polynomial-countermodel shifts are integral & \lrefx{Erdos251/PrimeGapDyadicTail.lean}{1654}{polynomialTailOrbit_shift_integral}\\
polynomial countermodel dyadic value & \lrefx{Erdos251/PolynomialGapSeriesValue.lean}{92}{tsum_polynomialGapDyadicTerm_eq}\\
polynomial countermodel rationality & \lrefx{Erdos251/PolynomialGapSeriesValue.lean}{109}{not_irrational_tsum_polynomialGapDyadicTerm}\\
real tail shifted-gap identity & \lrefx{Erdos251/RealPrimeGapTail.lean}{31}{realPrimeGapTail_eq_tsum_shifted_gaps}\\
real tail recurrence & \lrefx{Erdos251/RealPrimeGapTail.lean}{48}{realPrimeGapTail_recurrence}\\
actual escape equivalence & \lrefx{Erdos251/RealPrimeGapTail.lean}{73}{irrational_primeSeries_iff_realPrimeGapTail_escape}\\
real adjacent-pair consumer & \lrefx{Erdos251/RealPrimeGapTail.lean}{99}{realTailShift_not_both_integral_of_small_pair_of_digit_ne}\\
real small-mismatch criterion & \lrefx{Erdos251/RealPrimeGapTail.lean}{120}{irrational_primeSeries_of_realPrimeGapTail_small_mismatch}\\
free-pair lattice & \lrefx{Erdos251/FreePairReduction.lean}{92}{free_pair_integral_iff_modEq}\\
free-pair lattice corollary & \lrefx{Erdos251/FreePairReduction.lean}{105}{exists_free_pair_lattice}\\
actual real tail recurrence & \lrefx{Erdos251/FreePairReduction.lean}{186}{primeGapRealTail_recurrence}\\
free-pair equivalence for the actual series & \lrefx{Erdos251/FreePairReduction.lean}{236}{irrational_primeGap_tsum_iff_cofinalFreePairNonintegral}\\
kernel-decided certificate & \lrefx{Erdos251/KernelDenominatorFloor.lean}{310}{cert_10000}\\
denominator floor, prime series & \lrefx{Erdos251/KernelDenominatorFloor.lean}{316}{kernel_denominator_floor}\\
denominator floor, gap series & \lrefx{Erdos251/KernelDenominatorFloor.lean}{324}{kernel_denominator_floor_primeGap}\\
bounded-perturbation construction & \lrefx{Erdos251/BoundedPerturbationCountermodel.lean}{193}{exists_bounded_perturbation}\\
binary-digit encoding & \lrefx{Erdos251/BoundedPerturbationCountermodel.lean}{134}{exists_bounded_perturbation_digits}\\
prime-anchored rational perturbation & \lrefx{Erdos251/BoundedPerturbationCountermodel.lean}{208}{exists_rational_bounded_perturbation_primeGap}\\
perturbed prime-gap series is rational & \lrefx{Erdos251/BoundedPerturbationCountermodel.lean}{224}{not_irrational_bounded_perturbation_primeGap}\\
perturbed gap bounds & \lrefx{Erdos251/BoundedPerturbationCountermodel.lean}{234}{perturbed_gap_bounds}\\
signed two-window normal form & \lrefx{Erdos251/AffineShiftEscape.lean}{113}{adjacent_small_mismatch_iff_signed_two_window}\\
pointwise affine identity & \lrefx{Erdos251/AffineCylinderCollapse.lean}{91}{ratAffinePowTwo_iff_evenIntegral}\\
cofinal affine circularity & \lrefx{Erdos251/AffineCylinderCollapse.lean}{162}{cofinal_affinePowTwo_escape_iff_not_eventuallyIntegral}\\
fixed-lattice circularity & \lrefx{Erdos251/AffineCylinderCollapse.lean}{308}{cofinal_blockResidue_escape_iff_not_eventuallyIntegral}\\
lcm-diagonal criterion & \lrefx{Erdos251/OrderLatticeDiagonal.lean}{153}{irrational_initial_iff_all_lcmDiagonal_nonintegral}\\
\end{tabular}
\fi
